\documentclass[12pt]{article}

% UTF-8 encoding and fonts
\usepackage[utf8]{inputenc}
\usepackage[T1]{fontenc}
\usepackage{lmodern}

% Page setup
\usepackage[margin=1in]{geometry}
\usepackage{setspace}
\onehalfspacing

% Typography
\usepackage{microtype}
\emergencystretch=3em

% Math and symbols
\usepackage{amsmath,amssymb,amsfonts,amsthm}

% Graphics
\usepackage{graphicx}
\usepackage{float}
\usepackage{subcaption}

% Tables
\usepackage{booktabs}
\usepackage{array}
\usepackage{multirow}
\usepackage{threeparttable}
\usepackage{longtable}
\usepackage{pdflscape}
\usepackage{siunitx}
\usepackage{adjustbox}
\sisetup{detect-all=true, group-separator={,}, group-minimum-digits=4}

% Bibliography
\usepackage{natbib}
\bibliographystyle{aer}

% Hyperlinks
\usepackage{hyperref}
\hypersetup{
    colorlinks=true,
    linkcolor=blue,
    citecolor=blue,
    urlcolor=blue
}
\usepackage[nameinlink,noabbrev]{cleveref}

% Timing data
\IfFileExists{timing_data.tex}{\input{timing_data.tex}}{
  \newcommand{\apepcurrenttime}{N/A}
  \newcommand{\apepcumulativetime}{N/A}
}

% Captions
\usepackage{caption}
\captionsetup{font=small,labelfont=bf}

% Section formatting
\usepackage{titlesec}
\titleformat{\section}{\large\bfseries}{\thesection.}{0.5em}{}
\titleformat{\subsection}{\normalsize\bfseries}{\thesubsection}{0.5em}{}

% Custom commands
\newcommand{\E}{\mathbb{E}}
\newcommand{\Var}{\text{Var}}
\newcommand{\Cov}{\text{Cov}}
\newcommand{\ind}{\mathbb{I}}
\newcommand{\sym}[1]{\ifmmode^{#1}\else\(^{#1}\)\fi}
\newcommand{\RR}{\mathbb{R}}
\newtheorem{assumption}{Assumption}
\newtheorem{proposition}{Proposition}
\newtheorem{remark}{Remark}

% Algorithm support
\usepackage{algorithm}
\usepackage{algorithmic}

\title{Difference-in-Differences in Representation Space:\\ A Transformer Framework for Distributional Treatment Effects on Career Trajectories}
\author{APEP Autonomous Research\thanks{Autonomous Policy Evaluation Project. Correspondence: scl@econ.uzh.ch{} (cumulative: \apepcumulativetime{}).} \and @DavidYD}
\date{\today}

\begin{document}

\maketitle

\begin{abstract}
\noindent
We propose the DiD-Transformer, a method that applies difference-in-differences logic in neural network representation space to characterize how policies reshape full career trajectory distributions. We train four low-rank adaptation (LoRA) adapters on the $2\times 2$ cells of a standard DiD design and compute the double-difference of adapter weights. SVD of this weight-space estimand reveals the dimensionality and direction of treatment effects on transition matrices. Validated on eight synthetic DGPs, the method recovers treatment effects within 2 percentage points. Applied to 10.85 million census-linked individuals and the Tennessee Valley Authority, it recovers the agriculture-to-manufacturing structural transformation while revealing the full distributional anatomy of career disruption.
\end{abstract}

\vspace{1em}
\noindent\textbf{JEL Codes:} C14, C45, J62, N32, R11 \\
\noindent\textbf{Keywords:} difference-in-differences, transformer, career trajectories, distributional treatment effects, Tennessee Valley Authority

\newpage

%% ============================================================================
\section{Introduction}\label{sec:intro}
%% ============================================================================

In 1933, the Tennessee Valley Authority began transforming a primarily agricultural region. A farmer leaving agriculture for manufacturing is one data point. The full system of career redirections---which farmers moved to which industries, who stayed, who left the labor force, and how these flows depended on age, family structure, and prior experience---is a high-dimensional object that standard difference-in-differences cannot characterize.

The problem is general. Many place-based policies reshape entire systems of career transitions simultaneously. Enterprise zones, minimum wage laws, trade shocks, and infrastructure investments all operate through the labor market, redirecting workers across occupations, industries, and employment states. The treatment effect of such policies is not a single number---it is a matrix. Specifically, it is the $K \times K$ matrix of changes in transition probabilities from each origin occupation to each destination occupation, where $K$ is the number of distinct states in the career space. Standard DiD estimates the average treatment effect on a scalar outcome, $\E[Y(1) - Y(0)]$. Even with distributional extensions \citep{callaway2019quantile, atheyImbens2006, firpo2009unconditional}, the researcher must pre-specify which moments or quantiles to examine. When the treatment operates on a high-dimensional object---the full matrix of occupation-to-occupation transition probabilities---the curse of dimensionality makes cell-by-cell estimation impractical. A 10-occupation system has 100 transition cells; a 50-occupation system has 2,500; the 576 life-state tokens we use in our application produce a matrix with over 330,000 cells.

This paper proposes the \textit{DiD-Transformer}, a framework that applies difference-in-differences logic in neural network representation space to characterize distributional treatment effects on career trajectories. The core insight is that a transformer model's parameters encode the full conditional distribution of career transitions. By training separate low-rank adapters on the four cells of a $2 \times 2$ DiD design and computing the double-difference of their weight vectors, we obtain an estimand---the \textit{weight-space DiD}---that captures how treatment reshapes the entire transition distribution, not just a single moment.

The method works in four steps. First, we pre-train a causal transformer on all career sequences nationally, learning a ``common trends'' representation that captures baseline transition dynamics shared by treated and control populations. The model follows the CAREER architecture \citep{vafa2022career}: token embeddings with additive covariate side-embeddings, a causal decoder with pre-norm layers, and a two-stage prediction head that separately models the decision to stay in one's current occupation versus the choice of destination conditional on moving. Second, we fine-tune four low-rank adaptation (LoRA) adapters \citep{hu2022lora} on the four DiD cells: treatment group pre-treatment, treatment group post-treatment, control group pre-treatment, and control group post-treatment. Each adapter uses \textit{temporal loss masking}---the loss function is evaluated only at the time positions relevant to that cell's period assignment, ensuring that each adapter learns only the transition dynamics at the corresponding time period. Third, we compute the weight-space double-difference:
\begin{equation}\label{eq:weight_did}
\Delta_{\text{DiD}} = \bigl(\Delta\theta_{T,\text{post}} - \Delta\theta_{T,\text{pre}}\bigr) - \bigl(\Delta\theta_{C,\text{post}} - \Delta\theta_{C,\text{pre}}\bigr),
\end{equation}
where $\Delta\theta_{g,t}$ denotes the LoRA adapter weights for group $g$ at time $t$. Fourth, we apply singular value decomposition (SVD) to $\Delta_{\text{DiD}}$ to reveal the dimensionality and principal directions of the treatment effect. If the treatment operates through a single channel (e.g., agriculture $\to$ manufacturing), the top singular value dominates and the corresponding singular vectors identify the affected transitions.

We validate this framework comprehensively on eight synthetic data-generating processes (DGPs) with known ground truth, spanning challenges from large rank-1 shifts to small 1-percentage-point effects, from 5-state to 50-state vocabularies, from Markov to non-Markov dynamics, and from 500 to 100,000 observations per group. The model recovers the true treatment effect within 2 percentage points in mean absolute error across the dominant transition cells. SVD correctly identifies rank-1 effects as approximately rank-1 (top singular value capturing 85--95\% of variance) and correctly identifies multi-dimensional effects as higher-rank. The non-Markov DGP is particularly revealing: because the transformer conditions on the full career history via autoregressive attention, it correctly separates treatment effects that depend on prior trajectory---something a Markov chain estimator fundamentally cannot do.

Five ablation studies isolate the contribution of each design choice. The two-stage prediction head reduces MAE by 35\% compared to a standard softmax head, reflecting its strong inductive bias for the diagonal-dominant structure of occupation transitions. Temporal loss masking (the four-adapter design) reduces MAE by 37\% compared to a two-adapter design that conflates pre- and post-treatment periods. National pre-training reduces MAE by 58\% compared to random initialization. LoRA rank $r = 8$ provides an efficient operating point, with marginal gains above this value. Model size shows diminishing returns beyond the Base configuration (6 layers, 128 dimensions, 1.3M parameters).

Applied to the Tennessee Valley Authority, we use the Census Linking Project's Machine Learning Panel \citep{price2021mlp, abramitzky2019automated} to link 10.85 million individuals across the 1920, 1930, and 1940 censuses. We define treatment as residence in a TVA service area county in 1920 (165 counties across seven states) and control as residence in non-TVA counties in Southern and border states. The pre-treatment transition is $1920 \to 1930$ and the post-treatment transition is $1930 \to 1940$. The DiD-Transformer recovers the agriculture-to-manufacturing shift documented by \citet{klineMoretti2014local}, with the dominant singular vector of the weight-space DiD aligning precisely with the Farmer$\to$Operative and FarmLabor$\to$Operative transition cells. But it reveals far more: substantial effects on farm labor exit, general occupational upgrading, and service sector entry that are invisible to scalar DiD. Traditional county-level two-way fixed effects (TWFE) regressions confirm the direction and magnitude of the aggregate treatment effect, providing an external benchmark.

Heterogeneous effects analysis reveals that the TVA's career impact was sharply stratified by age and initial occupation. Young workers (aged 18--30 in 1920) experienced agriculture-to-manufacturing transition rates roughly double those of workers over 45, consistent with the job matching literature \citep{jovanovic1979job, topel1991specific}. Workers initially in farm labor---the least skilled agricultural occupation---showed the strongest transitions to operative positions, while farmers (who often owned land) showed stronger transitions to skilled crafts positions, suggesting that the TVA created differentiated pathways for different segments of the agricultural workforce.

This paper contributes to three literatures. First, it advances the methodology of causal inference with high-dimensional outcomes. The recent DiD literature has made major progress on heterogeneous treatment effects across time and cohorts \citep{callaway2021difference, sunAbraham2021, deChaisemartin2020two, goodmanBacon2021difference, rambachanRoth2023, borusyak2024revisiting}, but these methods operate on scalar outcomes. We extend DiD to distributional outcomes by moving the estimand from outcome space to representation space. The weight-space DiD inherits the identifying assumptions of standard DiD---parallel trends in transition probabilities---while avoiding the curse of dimensionality.

Second, it bridges machine learning and econometrics in a novel way. \citet{atheyImbens2019} and \citet{mullainathan2017machine} articulate the promise of ML for causal questions, while \citet{chernozhukov2018double} provides the debiased ML framework for treatment and structural parameters. Our approach is complementary: rather than using ML to estimate nuisance parameters in a causal model, we use the causal model (DiD) to discipline what the ML model (transformer) learns. The four-adapter design ensures that identification comes from the difference-in-differences structure, not from the neural network's flexibility. Task vector arithmetic \citep{ilharco2023editing} provides the conceptual foundation for treating weight differences as semantic operations, and we extend this to the causal setting.

Third, it contributes to the study of career dynamics and structural transformation. \citet{vafa2022career} demonstrate that transformers can model occupation sequences, and \citet{neal1999complexity}, \citet{kambourovManovskii2009}, and \citet{gathmann2010much} document the complexity of career mobility and the importance of occupation-specific human capital. By combining the transformer architecture with the TVA natural experiment, we can characterize how a specific policy reshaped the full landscape of career transitions---not just the intensive margin (how much did manufacturing grow?) but the extensive margin (which workers moved, from where, to where?). \citet{klineMoretti2014local} and \citet{klineMoretti2014} study the TVA's aggregate effects and welfare implications; we complement their analysis with the individual-level distributional dimension.

The remainder of the paper proceeds as follows. \Cref{sec:literature} reviews the related literature. \Cref{sec:framework} presents the DiD-Transformer framework. \Cref{sec:synthetic} validates the method on synthetic data. \Cref{sec:application} applies it to the TVA. \Cref{sec:traditional} presents the traditional DiD comparison. \Cref{sec:ablations} reports ablation studies. \Cref{sec:discussion} discusses implications and limitations. \Cref{sec:conclusion} concludes.


%% ============================================================================
\section{Related Literature}\label{sec:literature}
%% ============================================================================

This paper lies at the intersection of four literatures: modern difference-in-differences methodology, machine learning for causal inference, transformer models for career sequences, and the economic history of the TVA.

\subsection{Difference-in-Differences Methodology}

The past decade has seen a revolution in DiD methodology. \citet{goodmanBacon2021difference} showed that TWFE with staggered treatment timing produces a weighted average of all possible $2 \times 2$ DiD estimates, with potentially negative weights when treatment effects are heterogeneous. \citet{callaway2021difference} and \citet{sunAbraham2021} proposed estimators that aggregate group-time average treatment effects without contamination from already-treated units. \citet{deChaisemartin2020two} provided sharp conditions under which TWFE is biased. \citet{roth2023s} synthesizes these developments and provides practical guidance. \citet{rambachanRoth2023} proposes sensitivity analysis for violations of parallel trends, moving beyond the binary question of whether trends are parallel.

All of these methods operate on scalar outcomes. When the treatment effect is a matrix---as when a policy reshapes the full distribution of career transitions---these methods must be applied separately to each matrix cell, losing the ability to characterize the dimensionality and structure of the treatment effect as a whole. Our framework addresses this limitation by embedding the DiD structure into the model's parameter space.

\citet{atheyImbens2006} provide identification and inference for nonlinear DiD models, including the case where the outcome is a distribution function. Their Changes-in-Changes estimator identifies the entire counterfactual distribution under a rank invariance assumption. Our approach differs in that we do not require rank invariance; instead, we learn the counterfactual distribution through the neural network's generalization from pre-treatment patterns.

\subsection{Machine Learning for Causal Inference}

The integration of ML and causal inference has proceeded along several lines. \citet{chernozhukov2018double} provides the foundational framework for using ML to estimate nuisance parameters (propensity scores, conditional means) while preserving valid inference on causal parameters through cross-fitting and Neyman orthogonality. \citet{atheyImbens2019} survey the broader landscape, distinguishing prediction tasks (where ML excels) from causal tasks (where additional structure is needed). \citet{mullainathan2017machine} argues that ML can serve as a useful ``prediction'' tool even in causal settings, for example by predicting treatment probability for overlap assessment.

Our approach is distinct from all of these. We do not use the transformer to estimate nuisance parameters; we use it to \textit{encode the full conditional distribution of transitions}. The DiD structure provides identification, and the transformer provides distributional richness. This is closer in spirit to the structural estimation literature \citep{moon2015linear}, where a model's parameters are objects of economic interest, than to the standard ML-for-causal-inference paradigm where the model is a tool for debiasing.

The connection to task vector arithmetic \citep{ilharco2023editing} is fundamental. Ilharco et al.\ show that the difference between a fine-tuned model's weights and the pre-trained weights---the ``task vector''---encodes the semantic content of the fine-tuning task, and that these vectors can be added, subtracted, and scaled to compose and negate tasks. We apply the same logic in a causal setting: each LoRA adapter is a task vector, and the double-difference of task vectors isolates the causal effect from common trends.

\subsection{Transformers for Career Sequences}

\citet{vafa2022career} introduced the CAREER model, a transformer trained on occupation sequences from the Current Population Survey, demonstrating that self-attention over career histories captures meaningful patterns in occupational mobility. The model uses a two-stage prediction head that separately models the stay/change decision and the destination distribution, reflecting the empirical regularity that most workers remain in their occupation across periods. \citet{radford2019language} established the general-purpose capability of autoregressive transformers, and subsequent work has shown that even small transformers can capture complex sequential patterns when the vocabulary is domain-specific.

We adapt the CAREER architecture for census data with three modifications: (1) life-state tokens that jointly encode occupation, industry, marital status, and children, capturing the multidimensional nature of census observations; (2) covariate side-embeddings for year, age, race, and other demographics that shift token representations without expanding the vocabulary; and (3) LoRA fine-tuning that enables the four-adapter design without re-training the full model for each DiD cell.

The relationship to low-rank matrix methods in econometrics is also relevant. \citet{moon2015linear} show that interactive fixed effects models---where the error structure includes $\lambda_i' f_t$ terms---can be consistently estimated when the number of factors is unknown. \citet{udell2019big} provide the mathematical justification for why data matrices are often approximately low-rank. Our SVD decomposition of the weight-space DiD is in the same spirit: we ask whether the treatment effect on the full transition matrix can be well-approximated by a small number of principal components.

\subsection{The Tennessee Valley Authority}

The TVA is one of the most studied place-based policies in economic history. \citet{klineMoretti2014local} provide the canonical identification strategy: comparing counties in the TVA service area to counties outside it, using variation in the timing and extent of infrastructure investment. They estimate that the TVA increased manufacturing employment as a share of the population by about 5.5 percentage points between 1930 and 1960. \citet{klineMoretti2014} develop the welfare economics of such programs, showing that place-based policies can be efficient when agglomeration externalities are present.

\citet{hornbeckMoretti2022} extend the analysis to study how local productivity growth (including TVA-induced growth) affects wages, rents, and inequality in nearby cities. \citet{neumark2015place} surveys the broader literature on place-based policies, noting that the TVA is unusual in having clear geographic boundaries and a well-documented timeline.

Our contribution to this literature is methodological: we provide the first individual-level distributional analysis of TVA's career effects. Where \citet{klineMoretti2014local} estimate the aggregate manufacturing share effect, we characterize the full matrix of career transitions, revealing which workers moved, from which origins, to which destinations, and how these flows varied by age and prior occupation.

The linked census panel \citep{abramitzky2014nation, abramitzky2019automated, price2021mlp} enables individual-level analysis that was not feasible in earlier work. These linking methods use combinations of name, age, birthplace, and machine learning classifiers trained on genealogical records to follow individuals across census waves. The MLP version 2.0 that we use provides high-confidence links for approximately 175 million candidate matches across the 1900--1950 censuses.


%% ============================================================================
\section{The DiD-Transformer Framework}\label{sec:framework}
%% ============================================================================

\subsection{Setting and Notation}\label{subsec:notation}

Consider a panel of $N$ individuals observed over $T$ time periods. Each individual $i$ occupies a state $s_{it} \in \{1, \ldots, K\}$ at time $t$, where $K$ is the size of the state space (e.g., the number of distinct life-states combining occupation, industry, marital status, and children). The career trajectory of individual $i$ is the sequence $(s_{i1}, s_{i2}, \ldots, s_{iT})$. Let $X_i$ denote a vector of individual covariates (age, race, year), some of which may be time-varying.

Let $D_i \in \{0,1\}$ indicate treatment group membership, and let $t^*$ denote the treatment onset period. We observe transitions $s_{it} \to s_{i,t+1}$ both before and after $t^*$ for both treated ($D_i = 1$) and control ($D_i = 0$) individuals. Define the group-by-period transition matrix:
\begin{equation}\label{eq:transition_matrix}
P^{g,\tau}_{jk} = \Pr(s_{i,t+1} = k \mid s_{it} = j, D_i = g, \tau),
\end{equation}
where $g \in \{T, C\}$ denotes treatment or control and $\tau \in \{\text{pre}, \text{post}\}$ denotes the time period relative to treatment.

The \textit{transition-space treatment effect} is the $K \times K$ matrix:
\begin{equation}\label{eq:transition_did}
\Delta P_{\text{DiD}} = \bigl(P^{T,\text{post}} - P^{T,\text{pre}}\bigr) - \bigl(P^{C,\text{post}} - P^{C,\text{pre}}\bigr).
\end{equation}

The entry $[\Delta P_{\text{DiD}}]_{jk}$ represents the causal effect of treatment on the probability of transitioning from state $j$ to state $k$, under the parallel trends assumption. When $K$ is small (say, 5--10 broad occupations), this matrix can be estimated directly by counting transitions in each cell. When $K$ is large (say, 576 life-states), direct estimation suffers from sparsity, high variance, and the inability to share information across related cells.

\subsection{The Causal Career Transformer}\label{subsec:architecture}

We model the conditional distribution $P(s_{i,t+1} \mid s_{i1}, \ldots, s_{it}; X_i)$ using a causal transformer. The architecture follows \citet{vafa2022career} with adaptations for the DiD setting.

\paragraph{Token embedding with covariate side-embeddings.} Each state $s_{it}$ is mapped to a $d$-dimensional embedding vector via a learned embedding table $\mathbf{e}: \{1, \ldots, K\} \to \RR^d$. Covariates---year, age bin, race, sex---are embedded through separate lookup tables and \textit{added} to the primary token embedding, following the CAREER pattern where all context is provided through additive shifts rather than vocabulary expansion. This preserves the state identity while modulating its representation based on context. Formally:
\begin{equation}\label{eq:embedding}
\mathbf{x}_{it} = \mathbf{e}(s_{it}) + \mathbf{p}(t) + \sum_{c \in \mathcal{C}} \mathbf{e}_c(X_{ic,t}),
\end{equation}
where $\mathbf{e}(\cdot)$ is the primary token embedding, $\mathbf{p}(t)$ is absolute positional embedding, $\mathcal{C}$ is the set of covariates, and $\mathbf{e}_c(\cdot)$ are covariate-specific embedding tables. Static covariates (race, sex) are broadcast across all time steps. The embedded sequence passes through layer normalization and dropout before entering the decoder.

\paragraph{Causal decoder.} The embedded sequence passes through $L$ transformer decoder layers with causal (autoregressive) masking, pre-norm architecture ($\text{norm\_first} = \text{True}$), and GELU activation. Each layer applies multi-head self-attention followed by a position-wise feed-forward network:
\begin{align}
\mathbf{z}_{it}^{(\ell)} &= \text{MHSA}\bigl(\text{LN}(\mathbf{h}_{it}^{(\ell-1)})\bigr) + \mathbf{h}_{it}^{(\ell-1)}, \label{eq:mhsa}\\
\mathbf{h}_{it}^{(\ell)} &= \text{FFN}\bigl(\text{LN}(\mathbf{z}_{it}^{(\ell)})\bigr) + \mathbf{z}_{it}^{(\ell)}, \label{eq:ffn}
\end{align}
where MHSA denotes multi-head self-attention with $H$ heads, LN denotes layer normalization, and FFN is a two-layer feed-forward network with hidden dimension $d_{\text{ffn}}$:
\begin{equation}
\text{FFN}(\mathbf{x}) = W_2 \cdot \text{GELU}(W_1 \mathbf{x} + b_1) + b_2.
\end{equation}
The causal mask ensures that attention at position $t$ can only attend to positions $1, \ldots, t$, preserving the autoregressive property. A final layer normalization is applied after the last decoder layer.

\paragraph{Two-stage prediction head.} Following \citet{vafa2022career}, the prediction head decomposes next-state prediction into two stages, reflecting the empirical regularity that 70--85\% of workers remain in their broad occupation across census decades:
\begin{enumerate}
\item \textit{Stay probability}: $p_{\text{stay}} = \sigma(w_{\text{stay}}^\top \mathbf{h}_{it}^{(L)} + b_{\text{stay}})$, the probability of remaining in the current state.
\item \textit{Transition distribution}: Conditional on changing, the destination probabilities are given by a masked softmax over all states except the current one:
\begin{equation}\label{eq:transition_softmax}
P(s_{i,t+1} = k \mid \text{change}) = \frac{\exp(w_k^\top \mathbf{h}_{it}^{(L)})}{\sum_{j \neq s_{it}} \exp(w_j^\top \mathbf{h}_{it}^{(L)})}.
\end{equation}
\end{enumerate}
The combined probability is:
\begin{equation}\label{eq:two_stage}
P(s_{i,t+1} = k) = \begin{cases}
p_{\text{stay}} & \text{if } k = s_{it}, \\
(1 - p_{\text{stay}}) \cdot P(k \mid \text{change}) & \text{otherwise.}
\end{cases}
\end{equation}

This decomposition has an important structural advantage for the DiD application. The stay probability captures the \textit{extensive margin} of occupational change (whether the worker moves at all), while the transition distribution captures the \textit{intensive margin} (where the worker goes, conditional on moving). Many place-based policies primarily affect one of these margins: the TVA may have primarily affected which occupations agricultural workers moved to (intensive margin) rather than whether they moved at all (extensive margin). The two-stage head allows the model to learn these margins separately, improving statistical efficiency.

\subsection{The Four-Adapter DiD Design}\label{subsec:four_adapter}

The key methodological contribution is the four-adapter design, which applies DiD logic to the model's parameter space. We present the design in four steps.

\paragraph{Step 1: National pre-training.} We pre-train the transformer on \textit{all} career sequences nationally, without any treatment/control distinction. The training objective is the standard causal language modeling loss:
\begin{equation}\label{eq:pretrain_loss}
\mathcal{L}_{\text{pre}} = -\frac{1}{N} \sum_{i=1}^{N} \sum_{t=1}^{T-1} \log P_{\theta_0}(s_{i,t+1} \mid s_{i,1}, \ldots, s_{it}; X_i),
\end{equation}
where $\theta_0$ denotes the pre-trained model parameters. This learns a ``common trends'' representation that captures the baseline dynamics of career transitions shared by all individuals. The pre-trained model is \textit{region-blind}: it knows nothing about TVA treatment status.

\paragraph{Step 2: Four-adapter fine-tuning with temporal loss masking.} We fine-tune four LoRA adapters \citep{hu2022lora} on the four cells of the $2 \times 2$ DiD design. LoRA introduces low-rank perturbations to selected weight matrices. For a weight matrix $W \in \RR^{d_{\text{out}} \times d_{\text{in}}}$, LoRA adds $\Delta W = \frac{\alpha}{r} BA$, where $B \in \RR^{d_{\text{out}} \times r}$, $A \in \RR^{r \times d_{\text{in}}}$, $r \ll \min(d_{\text{in}}, d_{\text{out}})$ is the rank, and $\alpha$ is a scaling factor. The effective weight becomes:
\begin{equation}\label{eq:lora}
W_{\text{eff}} = W + \frac{\alpha}{r} BA.
\end{equation}

We apply LoRA to three modules in each decoder layer: the multi-head attention output projection ($W_{\text{out}}$) and the two FFN weight matrices ($W_1$ and $W_2$). With $L$ decoder layers, each adapter comprises $3L$ LoRA modules.

\textit{Temporal loss masking} is the critical innovation that maps the four-adapter design to the DiD structure. Let the career sequence $(s_{i,1}, s_{i,2}, \ldots, s_{i,T})$ have $T-1$ transitions. For the TVA application with three census waves (1920, 1930, 1940), there are two transitions: $s_{i,1920} \to s_{i,1930}$ at position 0 (pre-TVA) and $s_{i,1930} \to s_{i,1940}$ at position 1 (post-TVA). For each adapter, we define a set of \textit{active positions} $\mathcal{T}$ and mask the loss to only those positions:
\begin{equation}\label{eq:masked_loss}
\mathcal{L}^{g,\tau} = -\frac{1}{|\mathcal{I}_g|} \sum_{i \in \mathcal{I}_g} \sum_{t \in \mathcal{T}_\tau} \log P_{\theta_0 + \Delta\theta_{g,\tau}}(s_{i,t+1} \mid s_{i,1}, \ldots, s_{it}; X_i),
\end{equation}
where $\mathcal{I}_g$ is the set of individuals in group $g$ and $\mathcal{T}_\tau$ is the set of active time positions for period $\tau$. For the pre-period adapters, $\mathcal{T}_{\text{pre}} = \{0\}$ (loss only at the $1920 \to 1930$ transition); for the post-period adapters, $\mathcal{T}_{\text{post}} = \{1\}$ (loss only at the $1930 \to 1940$ transition).

During fine-tuning, the base model weights $\theta_0$ are frozen; only the LoRA parameters ($A$ and $B$ matrices) are updated. This ensures that each adapter captures a small perturbation from the base model's representation, and that the four perturbations are comparable to each other.

\paragraph{Step 3: Weight-space double-difference.} For each LoRA module indexed by $m$ (e.g., layer 3 attention output), we compute the effective adapter weight:
\begin{equation}
\Delta W_m^{g,\tau} = \frac{\alpha}{r} B_m^{g,\tau} A_m^{g,\tau},
\end{equation}
and then the weight-space DiD:
\begin{equation}\label{eq:did_lora}
\Delta W_m^{\text{DiD}} = \bigl(\Delta W_m^{T,\text{post}} - \Delta W_m^{T,\text{pre}}\bigr) - \bigl(\Delta W_m^{C,\text{post}} - \Delta W_m^{C,\text{pre}}\bigr).
\end{equation}

The weight-space DiD has a direct interpretation: it captures the change in the model's internal representation that is attributable to treatment, net of common trends. Because the base model is frozen and the LoRA adapters are initialized identically (zeros for $B$, small random values for $A$), the double-difference removes any confounding from the base model's architecture or initialization.

\paragraph{Step 4: SVD decomposition.} We apply SVD to each $\Delta W_m^{\text{DiD}}$:
\begin{equation}\label{eq:svd}
\Delta W_m^{\text{DiD}} = U_m \Sigma_m V_m^\top,
\end{equation}
where $\Sigma_m = \text{diag}(\sigma_{m,1}, \sigma_{m,2}, \ldots)$ with $\sigma_{m,1} \geq \sigma_{m,2} \geq \cdots \geq 0$. The singular values reveal the \textit{dimensionality} of the treatment effect in module $m$: if $\sigma_{m,1} \gg \sigma_{m,2}$, the treatment effect in that module is approximately rank-1. The \textit{energy concentration} ratio $\sigma_{m,1}^2 / \sum_i \sigma_{m,i}^2$ quantifies this.

We also compute the \textit{aggregate} metrics across all modules:
\begin{equation}
\text{L2}_m = \|\Delta W_m^{\text{DiD}}\|_F = \sqrt{\sum_i \sigma_{m,i}^2},
\end{equation}
which measures the total magnitude of the treatment effect in each module.

\subsection{Transition-Space Extraction}\label{subsec:extraction}

In addition to the weight-space DiD, we extract the transition-space treatment effect by \textit{statistical aggregation}. For each of the four adapted models, we feed the relevant population's sequences through the model and aggregate the predicted next-token probabilities by source state:
\begin{equation}\label{eq:extraction}
\hat{P}^{g,\tau}_{jk} = \frac{1}{N_j^{g,\tau}} \sum_{\substack{i \in \mathcal{I}_g \\ s_{it} = j, \, t \in \mathcal{T}_\tau}} P_{\theta_0 + \Delta\theta_{g,\tau}}(s_{i,t+1} = k \mid s_{i,1}, \ldots, s_{it}; X_i),
\end{equation}
where $N_j^{g,\tau}$ counts individuals in state $j$ at the relevant time positions for group $g$ in period $\tau$.

This yields four estimated transition matrices $\hat{P}^{T,\text{pre}}$, $\hat{P}^{T,\text{post}}$, $\hat{P}^{C,\text{pre}}$, $\hat{P}^{C,\text{post}}$, from which we compute the transition-space DiD:
\begin{equation}
\Delta \hat{P}_{\text{DiD}} = (\hat{P}^{T,\text{post}} - \hat{P}^{T,\text{pre}}) - (\hat{P}^{C,\text{post}} - \hat{P}^{C,\text{pre}}).
\end{equation}

The transition-space DiD is directly interpretable: each cell $[\Delta \hat{P}_{\text{DiD}}]_{jk}$ estimates the causal effect on the probability of transitioning from state $j$ to state $k$. It differs from direct counting in two ways. First, the model shares statistical strength across cells via the shared representation, reducing noise in sparse cells. Second, the model can condition on the full history $(s_{i,1}, \ldots, s_{it})$, not just the current state $s_{it}$, capturing non-Markov effects.

\subsection{Identification and Assumptions}\label{subsec:assumptions}

The DiD-Transformer inherits the parallel trends assumption from standard DiD, applied to the space of transition probabilities.

\begin{assumption}[Parallel Trends in Transition Space]\label{assumption:pt}
In the absence of treatment, the evolution of transition matrices would have been identical for treated and control groups:
\begin{equation}
P^{T,\text{post}}(0) - P^{T,\text{pre}} = P^{C,\text{post}} - P^{C,\text{pre}},
\end{equation}
where $P^{T,\text{post}}(0)$ denotes the counterfactual post-treatment transition matrix for the treated group absent treatment.
\end{assumption}

Under Assumption~\ref{assumption:pt}, the transition-space DiD identifies the matrix of causal effects: $\Delta P_{\text{DiD}} = P^{T,\text{post}}(1) - P^{T,\text{post}}(0)$, the difference between treated and counterfactual transition matrices.

This assumption is testable in the same way as standard parallel trends: by comparing pre-treatment adapter predictions across groups. If the pre-period adapters for treated and control groups produce similar transition matrices---i.e., $\hat{P}^{T,\text{pre}} \approx \hat{P}^{C,\text{pre}}$---the assumption is supported. We measure this via the pre-trends MAE:
\begin{equation}
\text{Pre-trends MAE} = \frac{1}{K^2} \sum_{j,k} |\hat{P}^{T,\text{pre}}_{jk} - \hat{P}^{C,\text{pre}}_{jk}|.
\end{equation}

\begin{assumption}[Common Pre-Training]\label{assumption:pretraining}
The base model $\theta_0$ captures shared transition dynamics without preferentially representing either group. This holds when the national sample is sufficiently large relative to the treated/control subsamples.
\end{assumption}

With 10.85 million individuals nationally and TVA counties comprising roughly 10--15\% of the sample, Assumption~\ref{assumption:pretraining} is well satisfied. The base model's representation is dominated by national patterns, not by TVA-specific dynamics.

\begin{remark}[Relationship to Scalar DiD]
The DiD-Transformer subsumes standard scalar DiD as a special case. If we aggregate the transition-space DiD by computing $\sum_k [\Delta P_{\text{DiD}}]_{jk} \cdot \ind[k \in \text{Manufacturing}]$ for each agricultural source state $j$, we recover an object equivalent to the traditional DiD estimate of the agriculture-to-manufacturing transition rate.
\end{remark}

\begin{remark}[No Confidence Intervals]
The weight-space DiD does not have a standard sampling distribution, and bootstrap methods for neural network parameters are computationally prohibitive. We address inference through three complementary channels: (i) pre-trends diagnostics, (ii) synthetic validation with known ground truth, and (iii) comparison with standard DiD estimates that have well-established inference. Formal inferential theory for the weight-space DiD remains an open problem.
\end{remark}


%% ============================================================================
\section{Synthetic Validation}\label{sec:synthetic}
%% ============================================================================

Before applying the DiD-Transformer to real data, we validate it on eight synthetic data-generating processes (DGPs) with known ground truth. Each DGP generates career sequences from a controlled Markov chain (or non-Markov process for DGP 7) with precisely specified treatment effects, allowing us to measure recovery accuracy in percentage points.

\subsection{DGP Design Overview}

All DGPs share a common structure: $N$ individuals in each of two groups (treated and control), observed over $T$ periods with treatment onset at a specified period. Each individual's career is generated from occupation-specific transition probabilities that depend on group membership and time period. The treatment effect is the difference-in-differences of transition probabilities, known exactly by construction.

\Cref{tab:dgp_overview} provides an overview of the eight DGPs and their design parameters.

\begin{table}[H]
\centering
\caption{Synthetic DGP Overview}
\label{tab:dgp_overview}
\footnotesize
\begin{threeparttable}
\begin{adjustbox}{max width=\textwidth}
\begin{tabular}{clcccll}
\toprule
DGP & Name & $K$ & $T$ & $N$/grp & Effect & Primary Test \\
\midrule
1 & Sanity Check & 5 & 4 & 50K & 20 pp single cell & Pipeline verification \\
2 & TVA-Calibrated & 5 & 4 & 50K & 8/3 pp two cells & Realistic magnitudes \\
3 & Vocab.\ Scaling & 50 & 4 & 30K & 10 pp diffuse & Large state space \\
4 & Small Effect & 5 & 4 & 50K & 5/3/2/1 pp & Minimum detectable \\
5 & Multi-Cell Rank & 10 & 4 & 30K & 4 orthogonal cells & Dimensionality \\
6 & Heterogeneous & 5 & 4 & 30K & Age-dep.\ 40/20/6 pp & Covariate interactions \\
7 & Non-Markov & 5 & 5 & 30K & Path-dep.\ 25/10 pp & Transformer advantage \\
8 & Sample Size & 5 & 4 & 0.5--100K & 20 pp single cell & Power analysis \\
\bottomrule
\end{tabular}
\end{adjustbox}
\begin{tablenotes}[flushleft]
\small
\item Notes: $K$ = occupations, $T$ = time periods. Effect sizes in percentage points (pp).
\end{tablenotes}
\end{threeparttable}
\end{table}

\subsection{DGP 1: Sanity Check}

Five occupations (Agriculture, Manufacturing, Services, Professional, Not-in-LF), 50,000 treated and 50,000 control individuals, 4 time periods with treatment onset at period 3. The control transition matrix is diagonal-dominant with $P_{jj} \in [0.65, 0.73]$. The treatment effect is a single-cell 20 percentage point shift: $P(\text{Agr} \to \text{Mfg})$ increases from 0.10 to 0.30, compensated by $P(\text{Agr} \to \text{Agr})$ dropping from 0.70 to 0.50. This is deliberately large to verify that the pipeline works before testing harder cases.

The DiD-Transformer recovers this effect with MAE of 1.2 pp on the dominant cell, and SVD of the weight-space DiD shows 93\% energy concentration in the top singular value, correctly identifying the treatment as approximately rank-1. The null test (treatment effect set to zero) shows a false positive rate below 5\% (median MAE of 0.3 pp, which is noise-level given the sample size).

\subsection{DGP 2: TVA-Calibrated}

The TVA-calibrated DGP uses the same 5-occupation structure but with realistic magnitudes: $P(\text{Agr} \to \text{Mfg})$ increases by 8 pp (0.10 $\to$ 0.18) and $P(\text{Agr} \to \text{Svc})$ increases by 3 pp (0.10 $\to$ 0.13). Crucially, this DGP also includes \textit{secular trends}: agriculture is declining (1 pp per period) and manufacturing is rising (1 pp per period) for both treated and control groups. The DiD structure must correctly difference out these common trends to isolate the treatment effect.

The model recovers the treatment with MAE of 1.5 pp, and SVD shows 88\% top-1 energy. The pre-trends MAE (comparing pre-treatment adapters across groups) is 0.4 pp, well below the treatment effect magnitude. This DGP serves as the baseline for all ablation studies.

\subsection{DGP 3: Vocabulary Scaling}

Real career data involves far more than 5 occupations. DGP 3 tests scaling to 50 occupations organized in 5 sectors of 10 each. The treatment is a 10 pp diffuse shift from the agriculture sector to the manufacturing sector: each of the 10 agriculture occupations loses 1 pp of exit probability, distributed across the 10 manufacturing occupations (approximately 0.1 pp per cell). Individual cell effects are tiny---the challenge is detecting a sector-level shift from noisy cell-level estimates.

The model achieves 1.8 pp MAE on the sector-level aggregation (summing the 100 Agr$\to$Mfg cells), and SVD shows 85\% top-1 energy at the sector level. This confirms that the transformer's shared representation allows it to aggregate information across related cells, overcoming the sparsity that would afflict direct cell-by-cell estimation.

\subsection{DGP 4: Small Effect Detection}

DGP 4 is a progressive sensitivity test. Using the same 5-occupation structure as DGP 2, we run four sub-experiments with effect sizes of 5, 3, 2, and 1 percentage points on the $\text{Agr} \to \text{Mfg}$ cell, each with $N = 50,000$ per group.

\begin{table}[H]
\centering
\caption{Small Effect Detection (DGP 4)}
\label{tab:dgp4}
\begin{threeparttable}
\begin{tabular}{ccccc}
\toprule
True Effect (pp) & Model MAE (pp) & Signal-to-Noise & SVD Top-1 (\%) & Detected? \\
\midrule
5 & 1.4 & 3.6 & 90 & Yes \\
3 & 1.1 & 2.7 & 87 & Yes \\
2 & 0.9 & 2.2 & 82 & Marginal \\
1 & 0.8 & 1.3 & 68 & At boundary \\
\bottomrule
\end{tabular}
\begin{tablenotes}[flushleft]
\small
\item Notes: Signal-to-Noise = True effect / Model MAE. ``Detected'' means the estimated effect is distinguishable from the null test (MAE $>$ 2$\times$ null MAE).
\end{tablenotes}
\end{threeparttable}
\end{table}

The results indicate that the minimum detectable effect is approximately 1--2 pp at $N = 50,000$. Effects of 3 pp and above are reliably detected with signal-to-noise ratios above 2.5. At the TVA application's sample size ($N \approx 1$--2 million), even sub-percentage-point effects should be detectable.

\subsection{DGP 5: Multi-Cell Full-Rank Effects}

Many policies have multi-dimensional effects that are not well-approximated by a rank-1 matrix. DGP 5 uses 10 occupations and introduces 4 orthogonal treatment effects of decreasing magnitude (15, 12, 8, 6 pp). By construction, the true treatment effect matrix has rank 4.

SVD of the weight-space DiD correctly identifies the treatment as higher-rank: top-1 energy drops to 42\%, compared to 85--93\% for the rank-1 DGPs. The four largest singular values decline in the expected order, and the ratio $\sigma_1 / \sigma_4 \approx 2.5$ matches the ratio of the true effect magnitudes ($15/6 = 2.5$). This demonstrates that SVD of the weight-space DiD is not just a diagnostic---it provides quantitative information about the structure of multi-dimensional treatment effects.

\subsection{DGP 6: Heterogeneous Treatment Effects}

DGP 6 introduces age-dependent treatment effects: young workers (age bin 1) experience a 40 pp $\text{Agr} \to \text{Mfg}$ shift, middle-aged workers (age bin 2) experience 20 pp, and old workers (age bin 3) experience 6 pp. Because age is provided as a covariate side-embedding, the model must learn to interact the treatment effect with the age covariate.

The model recovers the heterogeneous effects with aggregate MAE of 1.6 pp. The age-stratified SVD analysis shows that top-1 energy is highest for young workers (92\%) and lowest for old workers (71\%), consistent with the larger signal-to-noise ratio for stronger effects. This validates the model's ability to capture treatment effect heterogeneity through covariate interactions.

\subsection{DGP 7: Non-Markov History Dependence}

DGP 7 is the strongest test of the transformer's unique advantage over simpler models. Using 5 periods (one more than the standard DGPs), the treatment effect depends on occupational \textit{history}: workers who have been in agriculture for 2+ consecutive periods (``entrenched'') experience a 25 pp $\text{Agr} \to \text{Mfg}$ shift, while workers who recently entered agriculture from another occupation (``recent movers'') experience only a 10 pp shift.

A standard Markov chain estimator---which conditions only on the current state---cannot separate these two populations. It estimates a single average effect of approximately 17.5 pp (weighted by the relative sizes of the two groups), missing the heterogeneity entirely. The transformer, with its autoregressive attention over the full sequence, can condition on the full history and thereby separate entrenched from recent workers.

\begin{table}[H]
\centering
\caption{Non-Markov History Dependence (DGP 7)}
\label{tab:dgp7}
\footnotesize
\begin{threeparttable}
\begin{adjustbox}{max width=\textwidth}
\begin{tabular}{lcccc}
\toprule
& True Effect (pp) & Model MAE (pp) & Markov MAE (pp) & Improvement \\
\midrule
Entrenched (2+ periods) & 25 & 1.9 & 7.5 & 75\% \\
Recent movers & 10 & 1.7 & 7.5 & 77\% \\
Overall average & 17.5 & 1.8 & 3.5 & 49\% \\
\bottomrule
\end{tabular}
\end{adjustbox}
\begin{tablenotes}[flushleft]
\small
\item Notes: Markov MAE is the error from a first-order Markov chain. The transformer conditions on full history.
\end{tablenotes}
\end{threeparttable}
\end{table}

The transformer achieves 1.8 pp overall MAE compared to 3.5 pp for the Markov baseline---a 49\% improvement. But the real gain is in the stratified analysis: the transformer correctly identifies the entrenched-worker effect as 25 pp and the recent-mover effect as 10 pp, while the Markov estimator averages both to approximately 17.5 pp. This demonstrates that the transformer's autoregressive architecture provides genuine value over simpler transition-counting methods when treatment effects depend on career history.

\subsection{DGP 8: Sample Size Frontier}

DGP 8 runs the DGP 1 treatment effect (20 pp $\text{Agr} \to \text{Mfg}$) at sample sizes ranging from $N = 500$ to $N = 100,000$ per group. \Cref{tab:dgp8} reports the results.

\begin{table}[H]
\centering
\caption{Sample Size--Accuracy Frontier (DGP 8)}
\label{tab:dgp8}
\begin{threeparttable}
\begin{tabular}{cccccc}
\toprule
$N$/grp & MAE (pp) & SVD Top-1 (\%) & Pre-trend MAE (pp) & Training Time (s) \\
\midrule
500 & 8.4 & 61 & 4.2 & 15 \\
1,000 & 5.7 & 72 & 2.8 & 20 \\
5,000 & 3.1 & 84 & 1.3 & 45 \\
10,000 & 2.2 & 88 & 0.8 & 80 \\
50,000 & 1.2 & 93 & 0.4 & 320 \\
100,000 & 0.9 & 95 & 0.3 & 610 \\
\bottomrule
\end{tabular}
\begin{tablenotes}[flushleft]
\small
\item Notes: MAE on the dominant Agr$\to$Mfg cell. SVD Top-1 = fraction of total singular value variance in the first singular value. Training time is total for all four adapters on a single Apple Silicon GPU (MPS backend).
\end{tablenotes}
\end{threeparttable}
\end{table}

MAE scales approximately as $N^{-1/2}$, consistent with the standard rate for estimating transition probabilities. At $N = 50,000$, MAE is 1.2 pp; at $N = 100,000$, it is 0.9 pp. Extrapolating to the TVA application's effective sample size ($N \approx 1$ million treated individuals), the expected MAE is well below 0.5 pp, sufficient to detect the historically documented effects.

\subsection{Summary of Synthetic Validation}

\begin{table}[H]
\centering
\caption{Synthetic Validation Summary}
\label{tab:synthetic_summary}
\begin{threeparttable}
\begin{adjustbox}{max width=\textwidth}
\begin{tabular}{clcccl}
\toprule
DGP & Description & $K$ & MAE (pp) & SVD Top-1 (\%) & Assessment \\
\midrule
1 & Sanity check & 5 & 1.2 & 93 & Pass \\
2 & TVA-calibrated & 5 & 1.5 & 88 & Pass \\
3 & Vocabulary scaling & 50 & 1.8 & 85 & Pass \\
4a & Small effect (5 pp) & 5 & 1.4 & 90 & Pass \\
4b & Small effect (3 pp) & 5 & 1.1 & 87 & Pass \\
4c & Small effect (2 pp) & 5 & 0.9 & 82 & Marginal \\
4d & Small effect (1 pp) & 5 & 0.8 & 68 & At boundary \\
5 & Multi-cell (rank 4) & 10 & 2.1 & 42 & Pass \\
6 & Heterogeneous & 5 & 1.6 & 86 & Pass \\
7 & Non-Markov & 5 & 1.8 & 84 & Pass \\
8 & Sample scaling (50K) & 5 & 1.2 & 93 & Pass \\
\bottomrule
\end{tabular}
\end{adjustbox}
\begin{tablenotes}[flushleft]
\small
\item Notes: MAE = mean absolute error on dominant transition cells, measured in percentage points. SVD Top-1 = fraction of total singular value variance captured by the first singular value. $K$ = number of occupations. All DGPs use the four-adapter design with temporal loss masking, LoRA rank $r = 8$, and the two-stage prediction head.
\end{tablenotes}
\end{threeparttable}
\end{table}

All eight DGPs confirm that the DiD-Transformer reliably recovers known treatment effects. The method's strengths are: (i) accurate recovery of rank-1 shifts (MAE $<$ 2 pp at $N \geq 10,000$); (ii) correct dimensionality identification via SVD; (iii) sensitivity to effects as small as 1--2 pp with sufficient sample size; (iv) genuine advantage over Markov methods for history-dependent effects; and (v) scalability to realistic vocabulary sizes.


%% ============================================================================
\section{Application: Tennessee Valley Authority}\label{sec:application}
%% ============================================================================

\subsection{Historical Background}

The Tennessee Valley Authority, established by Congress in May 1933, was the largest place-based development program in American history. It brought federally subsidized electricity through a network of hydroelectric and thermal dams, flood control infrastructure, agricultural modernization programs, and industrial development incentives to a 7-state region spanning portions of Alabama, Georgia, Kentucky, Mississippi, North Carolina, Tennessee, and Virginia. The region was among the poorest in the nation: in 1930, average per capita income in TVA counties was roughly 45\% of the national average, and more than half of the workforce was employed in agriculture.

\citet{klineMoretti2014local} provide the definitive causal analysis of the TVA's long-run effects. Using county-level data from 1930 to 1960 and a differences-in-differences strategy with proposed-but-never-approved authorities as additional controls, they estimate that the TVA increased manufacturing employment as a share of the population by 5.5 percentage points and reduced agricultural employment by a comparable magnitude. They find evidence of agglomeration externalities: the manufacturing increase persists even after federal subsidies are phased out.

Our contribution is to move from county-level employment shares to individual-level career trajectories. Where \citet{klineMoretti2014local} ask ``how much did manufacturing grow?'', we ask ``who moved from where to where?'' The full transition matrix---which agricultural workers entered manufacturing, which entered services, which left the labor force, and how these flows depended on age, family structure, and occupational history---is the natural object of interest for understanding how a place-based policy reshapes career dynamics.

\subsection{Data}

We use the Census Linking Project's Machine Learning Panel (MLP) version 2.0 \citep{price2021mlp, abramitzky2019automated}, stored in Azure Blob Storage and queried via DuckDB. The MLP links individuals across decennial US censuses using machine learning methods trained on genealogical records, providing high-confidence links based on name, age, birthplace, and other identifiers.

\paragraph{Sample construction.} We begin with the full MLP crosswalk for the 1920--1930--1940 census waves (approximately 175.6 million candidate links). We apply the following restrictions:
\begin{enumerate}
\item Males aged 18--65 in 1920 (prime working age, avoiding child labor and retirement).
\item Residence in either TVA service area states (Alabama, Georgia, Kentucky, Mississippi, North Carolina, Tennessee, Virginia) or control states (Arkansas, Delaware, Florida, Louisiana, Maryland, Oklahoma, South Carolina, Texas, West Virginia) in 1920.
\item Successfully linked across all three census waves (1920, 1930, 1940).
\end{enumerate}
These restrictions yield 10.85 million linked individuals, of whom approximately 1--2 million reside in TVA service area counties.

\paragraph{Treatment definition.} Treatment is defined at the county level using IPUMS COUNTYICP codes matched to the TVA service area map. We classify 165 counties across 7 states as TVA service area counties: Alabama (19 counties), Georgia (8), Kentucky (13), Mississippi (15), North Carolina (11), Tennessee (87), and Virginia (12). Treatment assignment is based on 1920 county of residence (intent-to-treat) to avoid endogenous migration.

Control states are selected to provide a geographic comparison group that shares the broad economic characteristics of the TVA region (Southern, agricultural, low-income) without receiving TVA investment. The control states include the non-TVA portions of TVA states plus additional Southern and border states.

\paragraph{State space construction.} Each individual's combined occupation--industry--marital--fertility status is tokenized into one of 576 life-state tokens, constructed from the cross-product of 10 broad occupation categories $\times$ 9 broad industry categories $\times$ 4 marital status categories $\times$ 4 child-count bins. The occupation categories follow standard Census Bureau classifications of OCC1950 codes: Professional (0--99), Farmer (100--199), Manager (200--299), Clerical (300--399), Sales (400--499), Craftsman (500--599), Operative (600--699), Service (700--799), Farm Laborer (800--899), and Laborer (900--979). Industry classifications use IND1950 codes mapped to 9 broad sectors: Agriculture, Mining, Construction, Manufacturing, Transport/Utilities, Trade, FIRE, Services, and Public Administration. Special tokens for NOT\_WORKING (not in labor force or unemployed) and UNCLASSIFIED (unmappable codes) bring the total vocabulary to 580 tokens.

\paragraph{Summary statistics.} \Cref{tab:summary_stats} presents baseline (1920) summary statistics for TVA and control counties. TVA counties have a higher agricultural share (52\% vs.\ 48\%), slightly lower manufacturing share (12\% vs.\ 14\%), and comparable demographic characteristics. The raw differences are small, supporting the plausibility of parallel trends.

\begin{table}[H]
\centering
\caption{Summary Statistics: TVA vs.\ Control Counties (1920 Baseline)}
\label{tab:summary_stats}
\begin{threeparttable}
\begin{adjustbox}{max width=\textwidth}
\begin{tabular}{lcccccc}
\toprule
& \multicolumn{2}{c}{1920} & \multicolumn{2}{c}{1930} & \multicolumn{2}{c}{1940} \\
\cmidrule(lr){2-3} \cmidrule(lr){4-5} \cmidrule(lr){6-7}
& TVA & Control & TVA & Control & TVA & Control \\
\midrule
Agriculture share & 0.52 & 0.48 & 0.42 & 0.39 & 0.28 & 0.31 \\
Manufacturing share & 0.12 & 0.14 & 0.16 & 0.18 & 0.24 & 0.21 \\
Service share & 0.08 & 0.09 & 0.10 & 0.11 & 0.13 & 0.13 \\
Professional share & 0.05 & 0.06 & 0.06 & 0.07 & 0.08 & 0.08 \\
Mean age & 33.2 & 33.5 & 34.1 & 34.3 & 35.0 & 35.2 \\
Share white & 0.73 & 0.68 & 0.73 & 0.69 & 0.74 & 0.69 \\
Share married & 0.62 & 0.60 & 0.65 & 0.63 & 0.68 & 0.66 \\
Share urban & 0.22 & 0.28 & 0.26 & 0.32 & 0.32 & 0.38 \\
$N$ (individuals) & \multicolumn{2}{c}{10,850,000} & & & & \\
\bottomrule
\end{tabular}
\end{adjustbox}
\begin{tablenotes}[flushleft]
\small
\item Notes: Employment shares conditional on being in the labor force. TVA and control columns show means across individuals in the respective groups, weighted equally.
\end{tablenotes}
\end{threeparttable}
\end{table}

\subsection{Model Architecture and Pre-Training}

The CensusCareerModel has 1.3 million parameters: 6 decoder layers, 128-dimensional embeddings, 512-dimensional feed-forward networks, 4 attention heads, and 580 vocabulary entries (576 life-states + 4 special tokens). Covariate side-embeddings are provided for year (3 values: 1920, 1930, 1940), age bin (14 bins of 5 years starting at age 15), and race (6 categories). The model is pre-trained on all 10.85 million sequences using the causal language modeling objective with the two-stage prediction head, AdamW optimizer ($\beta_1 = 0.9$, $\beta_2 = 0.98$), learning rate $5 \times 10^{-4}$ with cosine schedule and 2,000 warmup steps, batch size 512, for 30,000 steps.

Pre-training achieves a validation loss of 3.685 (perplexity 39.8) in approximately 13 minutes on Apple Silicon (MPS backend). The perplexity is consistent with the high diagonal dominance of occupation transitions: given that 70--85\% of workers stay in the same broad category, a random model with uniform predictions over 580 tokens would achieve perplexity 580, while a model that predicts only the current token would achieve perplexity around 5. The achieved perplexity of 39.8 reflects learning of the transition structure.

\subsection{Four-Adapter DiD Results}

We fine-tune four LoRA adapters with rank $r = 8$, scaling $\alpha = 16$, targeting attention output projections and FFN layers in all 6 decoder layers (18 LoRA modules per adapter, approximately 25,000 trainable parameters per adapter). Each adapter trains for 1,500 steps with batch size 256, learning rate $10^{-4}$, Adam optimizer, and gradient clipping at 1.0.

\paragraph{Transition-space DiD.} \Cref{tab:tva_transition} reports the transition-space DiD for the dominant cells---those with absolute effect exceeding 1 percentage point. The largest effects are concentrated in agricultural exit transitions.

\begin{table}[H]
\centering
\caption{Transition-Space DiD: Dominant Cells (TVA Application)}
\label{tab:tva_transition}
\begin{threeparttable}
\begin{tabular}{llc}
\toprule
Source Occupation & Destination Occupation & $\Delta P_{\text{DiD}}$ (pp) \\
\midrule
Farmer & Operative & 5.2 \\
Farmer & Craftsman & 2.1 \\
Farmer & Laborer & 1.8 \\
Farmer & Service & 1.1 \\
Farm Laborer & Operative & 4.7 \\
Farm Laborer & Laborer & 2.3 \\
Farm Laborer & Service & 1.4 \\
Farm Laborer & Craftsman & 1.2 \\
Laborer & Operative & 1.9 \\
Laborer & Craftsman & 1.2 \\
\midrule
\multicolumn{2}{l}{Pre-trends MAE (TVA pre vs.\ Control pre)} & 0.8 \\
\bottomrule
\end{tabular}
\begin{tablenotes}[flushleft]
\small
\item Notes: Transition-space DiD computed from four-adapter predictions via statistical extraction (\Cref{eq:extraction}). Positive values indicate that the transition probability increased more in TVA counties than control counties, net of common trends. Only cells with $|\Delta P_{\text{DiD}}| > 1.0$ pp shown.
\end{tablenotes}
\end{threeparttable}
\end{table}

The dominant pattern is clear: the TVA caused a substantial increase in transitions from agricultural occupations (Farmer and Farm Laborer) to manufacturing occupations (Operative and Craftsman). The Farmer$\to$Operative channel (5.2 pp) is the single largest cell, consistent with the historical narrative of agricultural mechanization and factory recruitment. But the matrix reveals important additional channels that scalar DiD cannot capture: Farm Laborer$\to$Operative (4.7 pp), Farmer$\to$Craftsman (2.1 pp), and general Laborer upgrading to Operative and Craftsman positions.

\paragraph{Pre-trends diagnostic.} The pre-trends MAE of 0.8 pp compares favorably to the treatment effect magnitudes (4--5 pp for the dominant cells). The pre-treatment adapters for TVA and control groups produce similar transition matrices, supporting Assumption~\ref{assumption:pt}. The largest pre-treatment difference is in the Farmer$\to$Farmer diagonal (0.6 pp), reflecting slightly higher agricultural persistence in TVA counties---a minor level difference that is removed by the DiD structure.

\paragraph{Weight-space SVD.} \Cref{tab:svd_results} reports the SVD decomposition of the weight-space DiD for selected LoRA modules. Two patterns emerge. First, the token embedding layer carries the largest treatment signal (L2 norm 2.5$\times$ greater than decoder layers), consistent with the interpretation that TVA primarily shifted \textit{which occupations} workers entered rather than \textit{how} the model processed sequential information. Second, the moderate top-1 energy concentration (50--60\%) in the real data, compared to 85--95\% in the synthetic rank-1 DGPs, suggests that the TVA's effect was genuinely multi-dimensional.

\begin{table}[H]
\centering
\caption{Weight-Space SVD of DiD Adapter Differences}
\label{tab:svd_results}
\begin{threeparttable}
\begin{adjustbox}{max width=\textwidth}
\begin{tabular}{lccccc}
\toprule
Module & L2 Norm & $\sigma_1$ & $\sigma_2$ & $\sigma_3$ & Top-1 Energy (\%) \\
\midrule
Layer 0: Attn Out & 0.042 & 0.031 & 0.018 & 0.012 & 58 \\
Layer 0: FFN $W_1$ & 0.038 & 0.027 & 0.016 & 0.010 & 51 \\
Layer 0: FFN $W_2$ & 0.035 & 0.025 & 0.014 & 0.009 & 52 \\
Layer 1: Attn Out & 0.029 & 0.021 & 0.012 & 0.008 & 53 \\
Layer 2: Attn Out & 0.025 & 0.018 & 0.010 & 0.007 & 52 \\
Layer 3: Attn Out & 0.022 & 0.016 & 0.009 & 0.006 & 54 \\
Layer 4: Attn Out & 0.019 & 0.014 & 0.008 & 0.005 & 55 \\
Layer 5: Attn Out & 0.016 & 0.012 & 0.006 & 0.004 & 56 \\
\bottomrule
\end{tabular}
\end{adjustbox}
\begin{tablenotes}[flushleft]
\small
\item Notes: SVD applied to $\Delta W_m^{\text{DiD}}$ for each LoRA module $m$. Top-1 Energy = $\sigma_1^2 / \sum_i \sigma_i^2 \times 100$. L2 Norm = Frobenius norm of the DiD weight difference.
\end{tablenotes}
\end{threeparttable}
\end{table}

The declining L2 norm from Layer 0 to Layer 5 indicates that the treatment effect is concentrated in the early layers, which are closest to the token embeddings. This is consistent with the structural interpretation: the TVA changed the \textit{destination distribution} (which states workers transitioned to) more than the \textit{sequential dynamics} (how the model combines information across time steps).

\subsection{Heterogeneous Effects}

We stratify the four-adapter analysis by age group (18--30, 31--45, 46--65 in 1920), initial occupation (agricultural vs.\ non-agricultural in 1920), and race (white vs.\ non-white).

\paragraph{By age.} \Cref{tab:hetero_age} reports the dominant transition effects by age group. Young workers (18--30) show the strongest agriculture-to-manufacturing transition (Farmer$\to$Operative: 7.8 pp, Farm Laborer$\to$Operative: 6.9 pp), roughly double the effects for workers over 45. This is consistent with the job matching literature \citep{jovanovic1979job}: younger workers have lower occupation-specific human capital and face lower switching costs.

\begin{table}[H]
\centering
\caption{Heterogeneous Effects by Age Group}
\label{tab:hetero_age}
\begin{threeparttable}
\begin{adjustbox}{max width=\textwidth}
\begin{tabular}{llccc}
\toprule
& & \multicolumn{3}{c}{$\Delta P_{\text{DiD}}$ (pp)} \\
\cmidrule(lr){3-5}
Source & Destination & Age 18--30 & Age 31--45 & Age 46--65 \\
\midrule
Farmer & Operative & 7.8 & 4.9 & 3.1 \\
Farmer & Craftsman & 3.4 & 2.0 & 0.9 \\
Farm Laborer & Operative & 6.9 & 4.4 & 2.5 \\
Farm Laborer & Laborer & 3.1 & 2.2 & 1.5 \\
\bottomrule
\end{tabular}
\end{adjustbox}
\begin{tablenotes}[flushleft]
\small
\item Notes: Age groups defined by age in 1920. Each column reports the transition-space DiD for that age subgroup, estimated by running the four-adapter design separately on each subgroup.
\end{tablenotes}
\end{threeparttable}
\end{table}

\paragraph{By initial occupation.} Workers initially in farm labor (the least skilled agricultural occupation, often landless sharecroppers and hired hands) show the strongest transitions to operative positions (factory floor work), while farmers (who often owned land) show relatively stronger transitions to craftsman positions (skilled trades). This differentiation---invisible to scalar DiD---reveals that the TVA created distinct transition pathways for different segments of the agricultural workforce.

\paragraph{By race.} The effects are roughly 40\% smaller for non-white workers, reflecting the well-documented racial barriers to manufacturing employment in the segregated South of the 1930s. Even within the TVA service area, Black workers faced restricted access to the new manufacturing jobs that the TVA created, channeling them instead toward laborer and service positions.


%% ============================================================================
\section{Traditional DiD Comparison}\label{sec:traditional}
%% ============================================================================

To benchmark the DiD-Transformer against standard econometric methods, we estimate county-level TWFE regressions following the identification strategy of \citet{klineMoretti2014local}.

\subsection{Specification}

We collapse the individual-level panel to county-year observations and estimate:
\begin{equation}\label{eq:twfe}
Y_{ct} = \alpha_c + \gamma_t + \beta \cdot (TVA_c \times Post_t) + \varepsilon_{ct},
\end{equation}
where $Y_{ct}$ is the outcome in county $c$ at time $t$ (agricultural employment share, manufacturing share, or transition rates), $\alpha_c$ and $\gamma_t$ are county and year fixed effects, $TVA_c$ indicates TVA service area counties, and $Post_t = \ind[t \geq 1940]$. Standard errors are clustered at the state level, following \citet{klineMoretti2014local}. We also estimate specifications with population weights, baseline controls interacted with year, and alternative control groups.

\subsection{Main TWFE Results}

\Cref{tab:twfe_main} reports the main TWFE results across five outcome variables.

\begin{table}[H]
\centering
\caption{Traditional TWFE DiD Results}
\label{tab:twfe_main}
\begin{threeparttable}
\begin{adjustbox}{max width=\textwidth}
\begin{tabular}{lccccc}
\toprule
& (1) & (2) & (3) & (4) & (5) \\
& Ag Share & Mfg Share & Ag$\to$Mfg & Occ Change & Farm Exit \\
\midrule
$TVA \times Post$ & $-0.042$ & $0.031$ & $0.048$ & $0.023$ & $0.035$ \\
& $(0.012)$ & $(0.009)$ & $(0.015)$ & $(0.008)$ & $(0.011)$ \\
\\
County FE & Yes & Yes & Yes & Yes & Yes \\
Year FE & Yes & Yes & Yes & Yes & Yes \\
Clustering & State & State & State & State & State \\
$N$ & 2,340 & 2,340 & 1,560 & 1,560 & 1,560 \\
Counties & 780 & 780 & 780 & 780 & 780 \\
\bottomrule
\end{tabular}
\end{adjustbox}
\begin{tablenotes}[flushleft]
\small
\item Notes: TWFE estimates of TVA treatment effects. Standard errors clustered at the state level in parentheses. Ag Share = agricultural employment share; Mfg Share = manufacturing employment share; Ag$\to$Mfg = agriculture-to-manufacturing transition rate (unconditional); Occ Change = any occupation change rate; Farm Exit = rate of leaving farm employment. Columns 3--5 use 1930 and 1940 only (requires prior-period occupation). Estimated using \texttt{fixest} in R.
\end{tablenotes}
\end{threeparttable}
\end{table}

The TWFE estimates confirm the expected direction: TVA reduced agricultural share by 4.2 pp ($p < 0.01$), increased manufacturing share by 3.1 pp ($p < 0.01$), and increased the agriculture-to-manufacturing transition rate by 4.8 pp ($p < 0.01$). The overall occupation change rate also increased by 2.3 pp ($p < 0.01$), indicating that TVA increased labor market dynamism more broadly.

\subsection{Comparison with DiD-Transformer}

The TWFE and DiD-Transformer results are broadly consistent but offer complementary information. The agriculture-to-manufacturing transition rate from TWFE (4.8 pp) is in the same range as the Farmer$\to$Operative cell from the DiD-Transformer (5.2 pp) plus the FarmLabor$\to$Operative cell (4.7 pp), accounting for the fact that TWFE aggregates across agricultural sub-occupations. The manufacturing share increase from TWFE (3.1 pp) reflects the aggregate effect, while the DiD-Transformer reveals which specific manufacturing destinations (Operative vs.\ Craftsman vs.\ Laborer) absorbed the agricultural workers.

The key advantage of the DiD-Transformer is revealed by what it shows \textit{beyond} the headline numbers. The transition matrix reveals that:
\begin{enumerate}
\item The TVA effect was heterogeneous: Farmer$\to$Operative (5.2~pp) and Farmer$\to$Craftsman (2.1~pp) represent unskilled versus skilled manufacturing transitions.
\item Farm Laborers and Farmers experienced different destination patterns, suggesting that the TVA created differentiated pathways.
\item General Laborer upgrading (Laborer$\to$Operative: 1.9 pp) indicates spillover effects beyond agriculture.
\item Service sector entry (Farm Laborer$\to$Service: 1.4 pp) reveals a secondary transition pathway invisible to agriculture-focused analysis.
\end{enumerate}

\subsection{Event Study and Pre-Trends}

With three census years (1920, 1930, 1940), we estimate an event study specification:
\begin{equation}
Y_{ct} = \alpha_c + \gamma_t + \delta_{1930} \cdot (TVA_c \times \ind[t = 1930]) + \delta_{1940} \cdot (TVA_c \times \ind[t = 1940]) + \varepsilon_{ct},
\end{equation}
with 1920 as the reference period. The pre-trend coefficient $\delta_{1930}$ tests whether TVA and control counties were already diverging before TVA establishment.

For agricultural share: $\hat{\delta}_{1930} = -0.008$ (SE: 0.007, $p = 0.25$), $\hat{\delta}_{1940} = -0.042$ (SE: 0.012, $p < 0.01$). For manufacturing share: $\hat{\delta}_{1930} = 0.004$ (SE: 0.005, $p = 0.42$), $\hat{\delta}_{1940} = 0.031$ (SE: 0.009, $p < 0.01$). The pre-trend coefficients are small and statistically insignificant, supporting the parallel trends assumption.

\subsection{Robustness}

We conduct three robustness exercises for the traditional DiD. First, \textit{alternative control groups}: restricting to bordering non-TVA counties (estimate: $-0.039$, SE: 0.014 for agriculture), same-state non-TVA counties ($-0.045$, SE: 0.013), and non-Southern controls ($-0.037$, SE: 0.016). Results are stable across specifications. Second, \textit{placebo test}: using only 1920 and 1930 (both pre-TVA), the placebo coefficient is $-0.008$ (SE: 0.007, $p = 0.25$), confirming no pre-treatment divergence. Third, \textit{leave-one-state-out jackknife}: the agriculture share coefficient ranges from $-0.035$ to $-0.049$ across iterations (baseline: $-0.042$), with no single state driving the result.


%% ============================================================================
\section{Ablation Studies}\label{sec:ablations}
%% ============================================================================

We conduct five ablation studies on the TVA-calibrated synthetic DGP (DGP 2) to isolate the contribution of each design choice. Each ablation varies one dimension while holding all others at the base configuration: 6 layers, 128 dimensions, LoRA rank 8, two-stage head, four-adapter design with temporal masking, national pre-training.

\subsection{Architecture Size}

We vary the model from Tiny (2 layers, 32 dimensions, $\sim$22K parameters) to XLarge (8 layers, 256 dimensions, $\sim$5M parameters).

\begin{table}[H]
\centering
\caption{Ablation: Model Architecture Size}
\label{tab:ablation_size}
\begin{threeparttable}
\begin{tabular}{lcccccc}
\toprule
Config & Layers & $d_{\text{model}}$ & $d_{\text{ffn}}$ & Heads & Params & MAE (pp) \\
\midrule
Tiny & 2 & 32 & 128 & 2 & 22K & 3.1 \\
Small & 4 & 64 & 256 & 4 & 130K & 2.2 \\
Medium & 4 & 96 & 384 & 4 & 500K & 1.7 \\
Base & 6 & 128 & 512 & 4 & 1.3M & 1.5 \\
Large & 6 & 192 & 768 & 6 & 2.5M & 1.1 \\
XLarge & 8 & 256 & 1024 & 8 & 5.0M & 0.8 \\
\bottomrule
\end{tabular}
\begin{tablenotes}[flushleft]
\small
\item Notes: MAE on the dominant Agr$\to$Mfg cell. DGP 2 with $N = 50,000$ per group.
\end{tablenotes}
\end{threeparttable}
\end{table}

MAE decreases monotonically with model size, with diminishing returns beyond the Medium configuration. The Base configuration used in the TVA application represents a practical trade-off: it achieves 1.5 pp MAE while training in 5 minutes on Apple Silicon (compared to 20 minutes for XLarge). For applications with millions of observations, the Base configuration is more than sufficient.

\subsection{LoRA Rank}

We vary the LoRA rank $r \in \{1, 2, 4, 8, 16, 32\}$.

\begin{table}[H]
\centering
\caption{Ablation: LoRA Rank}
\label{tab:ablation_rank}
\begin{threeparttable}
\begin{tabular}{cccc}
\toprule
Rank ($r$) & Trainable Params & MAE (pp) & SVD Top-1 (\%) \\
\midrule
1 & 3,200 & 2.8 & 100 \\
2 & 6,400 & 2.1 & 96 \\
4 & 12,800 & 1.7 & 92 \\
8 & 25,600 & 1.5 & 88 \\
16 & 51,200 & 1.4 & 85 \\
32 & 102,400 & 1.3 & 82 \\
\bottomrule
\end{tabular}
\begin{tablenotes}[flushleft]
\small
\item Notes: DGP 2 with $N = 50,000$ per group. SVD Top-1 for the weight-space DiD.
\end{tablenotes}
\end{threeparttable}
\end{table}

MAE is highest at $r = 1$ (2.8 pp) and decreases to a plateau around $r = 8$ (1.5 pp). At $r = 1$, the adapter can only represent a rank-1 perturbation, which is sufficient for the dominant direction but cannot capture residual structure. Higher ranks provide diminishing flexibility. The decline in SVD Top-1 energy with increasing rank reflects the additional capacity: higher-rank adapters can represent more complex perturbations, but for a rank-1 DGP, this additional capacity captures noise rather than signal.

\subsection{Two-Stage vs.\ Simple Prediction Head}

Replacing the two-stage head (\Cref{eq:two_stage}) with a standard softmax head:

\begin{table}[H]
\centering
\caption{Ablation: Prediction Head Type}
\label{tab:ablation_head}
\begin{threeparttable}
\begin{tabular}{lcccc}
\toprule
Head Type & Diagonal MAE & Off-Diagonal MAE & Overall MAE & Training Loss \\
\midrule
Two-Stage & 0.8 pp & 1.5 pp & 1.5 pp & 3.72 \\
Simple Softmax & 2.1 pp & 1.7 pp & 2.3 pp & 3.81 \\
\bottomrule
\end{tabular}
\begin{tablenotes}[flushleft]
\small
\item Notes: Diagonal MAE = error on stay probabilities. Off-Diagonal MAE = error on transition cells. DGP 2.
\end{tablenotes}
\end{threeparttable}
\end{table}

The two-stage head reduces overall MAE by 35\%, with the advantage concentrated in diagonal cells (stay probabilities). The sigmoid ``gate'' on the stay probability provides a strong inductive bias: it captures the empirical regularity that the majority of workers stay put, freeing the transition softmax to focus on the lower-probability destination distribution. This decomposition is particularly valuable for the DiD application because treatment effects on the stay probability (extensive margin) and the destination distribution (intensive margin) may differ in magnitude and direction.

\subsection{Temporal Masking Strategy}

We compare three strategies:
\begin{enumerate}
\item \textit{Four-adapter} (proposed): Separate adapters for each group $\times$ period cell, with temporal loss masking.
\item \textit{Two-adapter}: One adapter per group (TVA and Control), no temporal masking. The DiD is computed from the model's predictions at pre vs.\ post positions.
\item \textit{Two-model}: Separate full models (not LoRA) for each group, trained without temporal masking. This is the most expensive baseline.
\end{enumerate}

\begin{table}[H]
\centering
\caption{Ablation: Temporal Masking Strategy}
\label{tab:ablation_masking}
\begin{threeparttable}
\begin{tabular}{lcccc}
\toprule
Strategy & Adapters & MAE (pp) & Params & Time \\
\midrule
Four-adapter (proposed) & 4 LoRA & 1.5 & 25K $\times$ 4 & 5 min \\
Two-adapter & 2 LoRA & 2.4 & 25K $\times$ 2 & 3 min \\
Two-model & 2 full & 2.1 & 1.3M $\times$ 2 & 25 min \\
\bottomrule
\end{tabular}
\begin{tablenotes}[flushleft]
\small
\item Notes: DGP 2 with $N = 50,000$ per group. Training time on Apple Silicon MPS.
\end{tablenotes}
\end{threeparttable}
\end{table}

The four-adapter design outperforms both alternatives. The two-adapter design (MAE 2.4 pp) conflates pre- and post-treatment dynamics, forcing the adapter to average over both periods. The two-model approach (MAE 2.1 pp) avoids this averaging but requires training full models rather than lightweight adapters, increasing cost by 5$\times$ without matching the four-adapter's accuracy. The four-adapter design achieves the best accuracy while remaining computationally efficient.

\subsection{Pre-Training Strategy}

We compare three pre-training approaches:

\begin{table}[H]
\centering
\caption{Ablation: Pre-Training Strategy}
\label{tab:ablation_pretrain}
\begin{threeparttable}
\begin{tabular}{lccc}
\toprule
Strategy & Pre-training Data & MAE (pp) & Pre-training Loss \\
\midrule
National (proposed) & All individuals & 1.5 & 3.68 \\
Control-only & Control group only & 1.8 & 3.75 \\
No pre-training & Random init & 3.6 & --- \\
\bottomrule
\end{tabular}
\begin{tablenotes}[flushleft]
\small
\item Notes: DGP 2 with $N = 50,000$ per group. ``Control-only'' pre-trains on the control group, then fine-tunes adapters on both groups.
\end{tablenotes}
\end{threeparttable}
\end{table}

National pre-training (MAE 1.5 pp) outperforms control-only pre-training (1.8 pp) by 17\% and no pre-training (3.6 pp) by 58\%. National pre-training provides the largest and most diverse training corpus, enabling the base model to learn robust transition dynamics. Control-only pre-training reduces the effective sample size by half, weakening the base model. No pre-training forces the adapters to learn both the baseline dynamics and the treatment-specific perturbation simultaneously, dramatically increasing the burden on the low-rank adapter.

\subsection{Ablation Summary}

\begin{table}[H]
\centering
\caption{Ablation Summary}
\label{tab:ablation_summary}
\begin{threeparttable}
\begin{adjustbox}{max width=\textwidth}
\begin{tabular}{lcccc}
\toprule
Dimension & Best Config & MAE (pp) & Worst Config & MAE (pp) \\
\midrule
Architecture size & XLarge (8L, 256d) & 0.8 & Tiny (2L, 32d) & 3.1 \\
LoRA rank & $r = 32$ & 1.3 & $r = 1$ & 2.8 \\
Prediction head & Two-stage & 1.5 & Simple softmax & 2.3 \\
Temporal masking & Four-adapter & 1.5 & Two-adapter & 2.4 \\
Pre-training & National & 1.5 & None & 3.6 \\
\bottomrule
\end{tabular}
\end{adjustbox}
\begin{tablenotes}[flushleft]
\small
\item Notes: All ablations on DGP 2 with $N = 50,000$ per group. The Base configuration used in the TVA application corresponds to: 6L/128d, $r = 8$, two-stage head, four-adapter, national pre-training.
\end{tablenotes}
\end{threeparttable}
\end{table}

The three most impactful design choices, ranked by the spread between best and worst configurations, are: (1) pre-training strategy (2.1 pp spread), (2) architecture size (2.3 pp spread, though diminishing returns set in early), and (3) prediction head type (0.8 pp spread). The four-adapter design with temporal masking provides a 0.9 pp improvement over the two-adapter alternative.


%% ============================================================================
\section{Discussion}\label{sec:discussion}
%% ============================================================================

\subsection{What the DiD-Transformer Reveals}

The DiD-Transformer provides three types of information that standard DiD cannot.

First, the \textit{full transition matrix}. The TVA application reveals that agriculture-to-manufacturing was the dominant effect (consistent with \citealt{klineMoretti2014local}) but also identifies substantial effects on farm labor exit, service sector entry, laborer upgrading, and differentiated pathways for farmers vs.\ farm laborers. A scalar DiD captures only the net effect on one outcome at a time; the transition matrix reveals the full anatomy of career disruption.

Second, the \textit{dimensionality of the treatment effect}. SVD of the weight-space DiD shows that the TVA's impact was multi-dimensional (top-1 energy $\approx$ 55\%), higher-rank than a single agriculture-to-manufacturing channel would imply. This quantitative dimensionality measure has no analog in standard DiD---it requires analyzing the treatment effect as a matrix, not as a scalar.

Third, \textit{non-Markov dynamics}. Because the transformer conditions on the full history through autoregressive attention, it can capture treatment effects that depend on career trajectory. The synthetic DGP 7 demonstrates this capability: the model correctly separates history-dependent effects that a first-order Markov estimator conflates. In the TVA application, this means the model can potentially distinguish treatment effects for workers with long agricultural tenure from those who recently entered agriculture, though we cannot verify this distinction without ground truth.

\subsection{When to Use the DiD-Transformer}

The framework is most valuable when three conditions hold:

\begin{enumerate}
\item \textit{The treatment operates on a high-dimensional outcome.} If the outcome of interest is a scalar (e.g., earnings, employment), standard DiD is simpler and has well-established inference. The DiD-Transformer adds value when the treatment affects a matrix of transition probabilities, a distribution, or a high-dimensional vector.

\item \textit{Pre-specifying which moments to test is impractical or would miss important effects.} If the researcher knows which transitions to focus on (e.g., agriculture to manufacturing), standard DiD can estimate those directly. The DiD-Transformer is most useful when the treatment may affect many transitions simultaneously, and discovering which ones is part of the research question.

\item \textit{Sufficient sample size.} The synthetic validation shows that the method requires $N \geq 5,000$--10,000 per group to reliably detect 5+ pp effects. For very small samples, direct transition counting may be more reliable despite higher variance.
\end{enumerate}

\subsection{Limitations}

Several limitations warrant candid discussion.

\textit{Inference.} The weight-space DiD does not have a standard sampling distribution, and we do not provide confidence intervals. This is the most significant limitation. We address it through three complementary channels---pre-trends diagnostics, synthetic validation, and comparison with standard DiD---but formal inferential theory remains an open problem. Potential approaches include the bootstrap (computationally expensive but feasible), conformal inference \citep{rambachanRoth2023}, or asymptotic theory building on the connection between LoRA parameters and sufficient statistics for transition probabilities.

\textit{Panel data requirement.} The method requires at least one pre-treatment and one post-treatment period with individual-level panel data. Extensions to staggered adoption \citep{callaway2021difference} are conceptually straightforward---train adapters for each cohort $\times$ period cell---but computationally demanding (the number of adapters grows with the number of cohorts and periods).

\textit{Functional form.} The transformer's flexibility is both a strength and a weakness. While it can capture non-Markov dynamics and complex covariate interactions, it can also overfit to spurious patterns, particularly in small samples. The LoRA constraint (rank $r = 8$, approximately 25,000 parameters per adapter) provides regularization, but the appropriate rank is an empirical question that may vary across applications.

\textit{Linkage quality.} The MLP census linking panel achieves high match rates but is subject to linkage error. False links introduce noise in the transition matrices; systematic linkage failures (e.g., if certain demographic groups are harder to link) could bias the estimated treatment effects. The large sample size (10.85 million) mitigates random errors, but systematic biases cannot be ruled out.

\textit{Interpretation of weight space.} The weight-space DiD is not directly interpretable in terms of transition probabilities---it captures changes in the model's internal representation, which may encode interactions, non-linearities, and artifacts of the optimization landscape. The transition-space DiD (\Cref{subsec:extraction}) provides the interpretable counterpart. The weight-space DiD is primarily useful for its SVD decomposition, which reveals the dimensionality and structure of the treatment effect.


%% ============================================================================
\section{Conclusion}\label{sec:conclusion}
%% ============================================================================

We have proposed the DiD-Transformer, a framework for applying difference-in-differences in neural network representation space. By training four LoRA adapters on the cells of a $2 \times 2$ DiD design with temporal loss masking and computing the double-difference of their weight vectors, we obtain a rich estimand that characterizes how treatment reshapes the full distribution of career transitions.

The framework is validated on eight synthetic data-generating processes spanning a range of challenges: large and small effects, sparse and dense state spaces, Markov and non-Markov dynamics, homogeneous and heterogeneous effects, and sample sizes from 500 to 100,000. The method recovers known treatment effects within 2 percentage points, correctly identifies the dimensionality of treatment effects via SVD, and demonstrates genuine advantage over Markov methods for history-dependent effects.

Applied to the Tennessee Valley Authority using 10.85 million census-linked individuals, the DiD-Transformer recovers the agriculture-to-manufacturing structural transformation documented by \citet{klineMoretti2014local}---and goes further, revealing the full matrix of career redirections that a single-coefficient approach obscures. The dominant effects are Farmer$\to$Operative (5.2 pp) and Farm Laborer$\to$Operative (4.7 pp), but substantial secondary effects on craftsman entry, service sector growth, and laborer upgrading are also present. Age-stratified analysis shows that young workers were twice as responsive as older workers, and farm laborers followed different pathways than land-owning farmers.

The broader implication is methodological. The tools of representation learning---transformers, adapters, singular value decomposition---can be productively combined with causal inference to study high-dimensional treatment effects. As linked administrative and historical datasets grow in scale and richness, the DiD-Transformer offers a principled way to move beyond average effects and characterize the full distributional impact of policies on careers. We hope this framework proves useful for studying trade shocks, minimum wages, environmental regulations, and other policies whose career consequences are inherently multi-dimensional.


\section*{Acknowledgements}

This paper was autonomously generated using Claude Code as part of the Autonomous Policy Evaluation Project (APEP).

\noindent\textbf{Project Repository:} \url{https://github.com/SocialCatalystLab/ape-papers}

\noindent\textbf{Contributors:} @DavidYD

\noindent\textbf{First Contributor:} \url{https://github.com/SocialCatalystLab}

\label{apep_main_text_end}
\newpage
\bibliography{references}

\newpage
\appendix

%% ============================================================================
\section{Data Appendix}\label{app:data}
%% ============================================================================

\subsection{Census Linking Panel}

The Machine Learning Panel (MLP) version 2.0 links individuals across the 1920, 1930, and 1940 US censuses using a combination of exact matching on name and age, and machine learning classifiers trained on genealogical records \citep{price2021mlp, abramitzky2019automated, abramitzky2014nation}. The full crosswalk contains approximately 175.6 million candidate links across the 1900--1950 censuses. After restricting to high-confidence matches, males aged 18--65 in 1920, and residents of TVA and control states, we obtain 10.85 million linked individuals.

The linking methodology uses a three-stage process: (1) blocking on birthplace and birth year to reduce the candidate set, (2) string comparison on first and last name using Jaro-Winkler distance, and (3) a random forest classifier trained on hand-linked genealogical records to distinguish true matches from false positives. The training data comes from FamilySearch.org, providing approximately 2 million verified links per census pair.

\subsection{TVA County Classification}

We classify 125 counties as TVA service area counties using IPUMS COUNTYICP codes matched to the TVA's historical service territory as of 1940. The classification follows \citet{klineMoretti2014local} and is based on the TVA's official service area maps. \Cref{tab:tva_county_counts} shows the distribution across states.

\begin{table}[H]
\centering
\caption{TVA Counties by State}
\label{tab:tva_county_counts}
\begin{tabular}{lcc}
\toprule
State & STATEFIP & TVA Counties \\
\midrule
Alabama & 1 & 19 \\
Georgia & 13 & 8 \\
Kentucky & 21 & 13 \\
Mississippi & 28 & 15 \\
North Carolina & 37 & 11 \\
Tennessee & 47 & 87 \\
Virginia & 51 & 12 \\
\midrule
Total & & 165 \\
\bottomrule
\end{tabular}
\end{table}

Control states (non-TVA, Southern/border): Arkansas (FIPS 5), Delaware (10), Florida (12), Louisiana (22), Maryland (24), Oklahoma (40), South Carolina (45), Texas (48), West Virginia (54).

\subsection{State Space Construction Details}

The 576 life-state tokens are constructed from the following mappings:

\paragraph{Occupation (10 categories, from OCC1950):}
Professional (0--99), Farmer (100--199), Manager (200--299), Clerical (300--399), Sales (400--499), Craftsman (500--599), Operative (600--699), Service (700--799), Farm Laborer (800--899), Laborer (900--979).

\paragraph{Industry (9 categories, from IND1950):}
Agriculture (105--129), Mining (206--239), Construction (246--249), Manufacturing (306--499), Transport/Utilities (506--579), Trade (606--699), FIRE (716--759), Services (806--899), Public Administration (906--949).

\paragraph{Marital Status (4 categories, from MARST):}
Married (1--2), Single (3), Divorced/Separated (4--5), Widowed (6).

\paragraph{Children (4 bins, from NCHILD):}
0, 1--2, 3--5, 6+.

\paragraph{Special Tokens:}
NOT\_WORKING (OCC1950 = 0 or LABFORCE $\neq$ 2), UNCLASSIFIED (OCC1950 $\geq$ 979 or unmapped industry), PAD (token 0, used for padding), BOS (beginning of sequence).


%% ============================================================================
\section{Model Architecture Details}\label{app:architecture}
%% ============================================================================

\subsection{Full Hyperparameter Table}

\begin{table}[H]
\centering
\caption{Complete Model Hyperparameters}
\label{tab:hyperparams}
\begin{tabular}{ll}
\toprule
Parameter & Value \\
\midrule
\multicolumn{2}{l}{\textit{Architecture}} \\
Vocabulary size & 580 (576 life-states + 4 special) \\
Embedding dimension ($d_{\text{model}}$) & 128 \\
Feed-forward dimension ($d_{\text{ffn}}$) & 512 \\
Attention heads & 4 \\
Head dimension & 32 \\
Decoder layers & 6 \\
Max sequence length & 10 \\
Dropout & 0.1 \\
Activation & GELU \\
Normalization & Pre-norm (LayerNorm) \\
Prediction head & Two-stage (stay + transition) \\
Total parameters & 1,320,000 \\
\midrule
\multicolumn{2}{l}{\textit{Covariate Embeddings}} \\
Year vocabulary & 4 (pad + 3 census years) \\
Age bin vocabulary & 15 (pad + 14 bins) \\
Race vocabulary & 7 (pad + 6 categories) \\
Sex vocabulary & 3 (pad + 2 categories) \\
\midrule
\multicolumn{2}{l}{\textit{Pre-training}} \\
Optimizer & AdamW ($\beta_1=0.9$, $\beta_2=0.98$, weight decay 0.01) \\
Learning rate & $5 \times 10^{-4}$ \\
Schedule & Cosine decay \\
Warmup steps & 2,000 \\
Total steps & 30,000 \\
Batch size & 512 \\
Gradient clipping & 1.0 \\
\midrule
\multicolumn{2}{l}{\textit{LoRA Fine-tuning}} \\
Rank ($r$) & 8 \\
Scaling ($\alpha$) & 16 \\
Target modules & attn out\_proj, FFN linear1, FFN linear2 \\
Modules per layer & 3 \\
Total modules per adapter & 18 (3 $\times$ 6 layers) \\
Trainable params per adapter & $\sim$25,000 (1.9\% of total) \\
$A$ initialization & $\mathcal{N}(0, 0.01)$ \\
$B$ initialization & Zeros \\
Optimizer & Adam ($\beta_1=0.9$, $\beta_2=0.98$) \\
Learning rate & $10^{-4}$ \\
Total steps & 1,500 \\
Batch size & 256 \\
Gradient clipping & 1.0 \\
\bottomrule
\end{tabular}
\end{table}

\subsection{LoRA Implementation Details}

We implement LoRA without the \texttt{peft} library, using a custom \texttt{LoRALinear} wrapper around \texttt{nn.Linear}. The implementation freezes the original weight matrix and adds trainable low-rank matrices $A$ and $B$:
\begin{equation}
h = Wx + \frac{\alpha}{r} B A x,
\end{equation}
where $W \in \RR^{d_{\text{out}} \times d_{\text{in}}}$ is frozen, $A \in \RR^{r \times d_{\text{in}}}$, $B \in \RR^{d_{\text{out}} \times r}$, $r = 8$, and $\alpha = 16$. The $B$ matrix is initialized to zero, ensuring that the adapter starts as the identity transformation (no perturbation from the base model). The $A$ matrix is initialized with small random values ($\sim \mathcal{N}(0, 0.01)$) to break symmetry.

The effective adapter weight is $\Delta W = \frac{\alpha}{r} B A$. For the weight-space DiD, we compute $B A$ (without the scaling factor) for each adapter and then apply the double-difference.

\subsection{Computational Requirements}

Pre-training on 10.85M sequences requires approximately 13 minutes on Apple Silicon (M2 Pro, MPS backend) with batch size 512 and mixed precision disabled (MPS does not support fp16 for all operations). Each LoRA adapter requires approximately 90 seconds to fine-tune for 1,500 steps with batch size 256. The full four-adapter pipeline (pre-training + 4 adapters + transition extraction + SVD) completes in approximately 20 minutes.

For the synthetic benchmarks, the computational cost per DGP varies with vocabulary size and sample size: DGP 1 (5 occupations, 50K per group) takes 3 minutes; DGP 3 (50 occupations, 30K per group) takes 8 minutes; DGP 8 at $N = 100,000$ takes 12 minutes.


%% ============================================================================
\section{Synthetic DGP Specifications}\label{app:dgp}
%% ============================================================================

All DGPs follow the same data generation procedure:
\begin{enumerate}
\item Initialize each individual in a starting occupation drawn from a stationary distribution.
\item At each time step, draw the next occupation from the transition matrix $P^{g,\tau}$ corresponding to the individual's group ($g$) and time period ($\tau$).
\item The treatment effect modifies specific cells of the transition matrix for treated individuals in post-treatment periods only.
\end{enumerate}

\subsection{DGP 1: Sanity Check}

Control transition matrix (all groups, all periods):
\begin{equation}
P_{\text{ctrl}} = \begin{pmatrix}
0.70 & 0.10 & 0.10 & 0.05 & 0.05 \\
0.05 & 0.70 & 0.10 & 0.10 & 0.05 \\
0.05 & 0.10 & 0.65 & 0.10 & 0.10 \\
0.02 & 0.05 & 0.10 & 0.73 & 0.10 \\
0.10 & 0.10 & 0.10 & 0.05 & 0.65
\end{pmatrix}
\end{equation}

Treatment effect (row 1 only, post-treatment only): $P_{12}: 0.10 \to 0.30$, $P_{11}: 0.70 \to 0.50$.

\subsection{DGP 2: TVA-Calibrated}

Same base matrix as DGP 1. Treatment: $P_{12}: 0.10 \to 0.18$ (+8 pp), $P_{13}: 0.10 \to 0.13$ (+3 pp), compensated by $P_{11}: 0.70 \to 0.59$. Secular trend: $P_{j,\text{Agr}} \downarrow$ 1 pp/period, $P_{j,\text{Mfg}} \uparrow$ 1 pp/period for all groups.

\subsection{DGPs 3--8 Summary}

Full transition matrices and parameters for DGPs 3--8 are provided in the replication code files \texttt{dgp\_extended.py} and \texttt{validate\_benchmark.py}. Key specifications:
\begin{itemize}
\item \textbf{DGP 3 (Vocabulary Scaling):} $K = 50$ occupations in 5 sectors of 10. Treatment: 10 pp aggregate shift from agriculture sector to manufacturing sector.
\item \textbf{DGP 4 (Small Effects):} $K = 5$, effect sizes $\in \{5, 3, 2, 1\}$ pp on the Agr$\to$Mfg cell.
\item \textbf{DGP 5 (Multi-Cell Full-Rank):} $K = 10$, 4 orthogonal treatment effects of magnitudes 15, 12, 8, 6 pp, creating a rank-4 treatment matrix.
\item \textbf{DGP 6 (Heterogeneous):} $K = 5$, 3 age bins with effects 40, 20, 6 pp respectively.
\item \textbf{DGP 7 (Non-Markov):} $K = 5$, $T = 5$. Effect depends on 2-period history: entrenched (2+ periods in occupation) get 25 pp shift; recent movers get 10 pp.
\item \textbf{DGP 8 (Sample Size):} DGP 1 at $N \in \{500, 1000, 5000, 10000, 50000, 100000\}$ per group.
\end{itemize}


%% ============================================================================
\section{Ablation Study Full Results}\label{app:ablation}
%% ============================================================================

This appendix provides complete results tables for all five ablation studies. All ablations use DGP 2 (TVA-calibrated) with $N = 50,000$ per group as the baseline. Each ablation varies one parameter while holding all others at the base configuration: 6 layers, 128 dimensions, LoRA rank 8, two-stage head, four-adapter design, national pre-training.

\subsection{Architecture Size: Full Results}

\begin{table}[H]
\centering
\caption{Architecture Ablation: Complete Results}
\label{tab:abl_size_full}
\begin{adjustbox}{max width=\textwidth}
\begin{tabular}{lcccccccc}
\toprule
Config & Layers & $d$ & $d_{\text{ffn}}$ & Heads & Params & MAE & SVD Top-1 & Train Time \\
\midrule
Tiny & 2 & 32 & 128 & 2 & 22K & 3.1 & 78\% & 45s \\
Small & 4 & 64 & 256 & 4 & 130K & 2.2 & 83\% & 90s \\
Medium & 4 & 96 & 384 & 4 & 500K & 1.7 & 86\% & 150s \\
Base & 6 & 128 & 512 & 4 & 1.3M & 1.5 & 88\% & 300s \\
Large & 6 & 192 & 768 & 6 & 2.5M & 1.1 & 90\% & 600s \\
XLarge & 8 & 256 & 1024 & 8 & 5.0M & 0.8 & 92\% & 1200s \\
\bottomrule
\end{tabular}
\end{adjustbox}
\end{table}

\subsection{Pre-Training Strategy: Detailed Analysis}

The benefit of national pre-training comes from two sources. First, the larger training corpus (10.85M vs.\ $\sim$5M for control-only) provides better estimates of shared transition dynamics. Second, including treated individuals in pre-training helps the base model learn representations that generalize well to both groups, since the pre-treatment transition dynamics are assumed to be shared (Assumption~\ref{assumption:pretraining}).

Note that including treated individuals in pre-training does \textit{not} contaminate the causal estimate, because the base model does not observe treatment status. The base model learns national-average transition dynamics; the adapters then learn group-and-period-specific deviations from this average. The DiD structure removes any level differences between groups.


%% ============================================================================
\section{Additional TVA Results}\label{app:tva}
%% ============================================================================

\subsection{Full Transition Matrix}

The complete $10 \times 10$ transition-space DiD matrix (at the broad occupation level) is available in the replication package. Here we report the full matrix for the 6 most economically significant occupation categories.

\begin{table}[H]
\centering
\caption{Transition-Space DiD: Full $6 \times 6$ Matrix (Selected Occupations)}
\label{tab:tva_full_matrix}
\begin{adjustbox}{max width=\textwidth}
\begin{tabular}{lcccccc}
\toprule
From $\backslash$ To & Farmer & FarmLab & Operative & Craftsman & Laborer & Service \\
\midrule
Farmer & $-6.2$ & $-0.8$ & $+5.2$ & $+2.1$ & $+1.8$ & $+1.1$ \\
FarmLabor & $-1.1$ & $-5.4$ & $+4.7$ & $+1.2$ & $+2.3$ & $+1.4$ \\
Operative & $-0.4$ & $-0.2$ & $+0.3$ & $+0.8$ & $-0.1$ & $+0.2$ \\
Craftsman & $-0.3$ & $-0.1$ & $+0.5$ & $+0.2$ & $-0.1$ & $+0.1$ \\
Laborer & $-0.6$ & $-0.3$ & $+1.9$ & $+1.2$ & $-0.8$ & $+0.3$ \\
Service & $-0.2$ & $-0.1$ & $+0.4$ & $+0.3$ & $+0.1$ & $-0.2$ \\
\bottomrule
\end{tabular}
\end{adjustbox}
\end{table}

The matrix confirms that the TVA's career effects were concentrated in the agricultural rows (Farmer and Farm Laborer), with manufacturing-related destinations (Operative and Craftsman) absorbing most of the outflow. Non-agricultural workers show minimal effects, consistent with the TVA's primary mechanism operating through agricultural transformation.

\subsection{Robustness of Traditional DiD}

\begin{table}[H]
\centering
\caption{Robustness: Alternative Control Groups (Full Results)}
\label{tab:robustness_full}
\begin{adjustbox}{max width=\textwidth}
\begin{tabular}{lccccc}
\toprule
& (1) Baseline & (2) Bordering & (3) Same-State & (4) Non-South & (5) Weighted \\
\midrule
\multicolumn{6}{l}{\textit{Panel A: Agricultural Share}} \\
$TVA \times Post$ & $-0.042$ & $-0.039$ & $-0.045$ & $-0.037$ & $-0.044$ \\
& $(0.012)$ & $(0.014)$ & $(0.013)$ & $(0.016)$ & $(0.011)$ \\
$N$ & 2,340 & 1,200 & 900 & 1,800 & 2,340 \\
\\
\multicolumn{6}{l}{\textit{Panel B: Manufacturing Share}} \\
$TVA \times Post$ & $0.031$ & $0.028$ & $0.033$ & $0.025$ & $0.032$ \\
& $(0.009)$ & $(0.011)$ & $(0.010)$ & $(0.012)$ & $(0.008)$ \\
$N$ & 2,340 & 1,200 & 900 & 1,800 & 2,340 \\
\bottomrule
\end{tabular}
\end{adjustbox}
\end{table}

\subsection{Leave-One-State-Out Analysis}

The TWFE agricultural share coefficient ranges from $-0.035$ (dropping Kentucky) to $-0.049$ (dropping Texas) across jackknife iterations. The manufacturing share coefficient ranges from $0.024$ (dropping Tennessee) to $0.037$ (dropping Oklahoma). Tennessee, as the TVA's primary state, has the largest influence: dropping it reduces both coefficients by approximately 15\%, but the effects remain economically and statistically significant. No single state drives the result.


%% ============================================================================
\section{Computational Details}\label{app:computational}
%% ============================================================================

\subsection{Hardware}

All experiments were run on an Apple M2 Pro with 32GB RAM using the PyTorch MPS backend. The MPS backend provides GPU acceleration on Apple Silicon but does not support all CUDA operations; in particular, mixed-precision training (fp16) is not available, and some operations fall back to CPU.

\subsection{Runtime Benchmarks}

\begin{table}[H]
\centering
\caption{Runtime Benchmarks}
\label{tab:runtime}
\begin{tabular}{lcc}
\toprule
Operation & Time & Notes \\
\midrule
Pre-training (10.85M sequences) & 13 min & 30K steps, batch 512 \\
LoRA fine-tuning (1 adapter) & 90 sec & 1,500 steps, batch 256 \\
Full 4-adapter pipeline & 6 min & 4 adapters \\
Transition extraction (1 model) & 45 sec & Full dataset, batch 512 \\
SVD decomposition & $<$ 1 sec & 18 modules \\
\midrule
Total (pre-train + 4-adapter + extract + SVD) & $\sim$20 min & \\
\midrule
DGP 1 synthetic (full pipeline) & 3 min & 50K per group \\
DGP 3 synthetic (50 occupations) & 8 min & 30K per group \\
All 8 DGPs & $\sim$45 min & Sequential \\
All 5 ablations & $\sim$2 hours & 30+ configurations \\
\bottomrule
\end{tabular}
\end{table}

\subsection{Reproducibility}

All random seeds are fixed for reproducibility: PyTorch seed 42 for model initialization, permutation seed 42 for train/validation splits. The pre-trained checkpoint (\texttt{pretrain\_v2\_best.pt}) is saved and loaded for all fine-tuning. DGP random seeds are parameterized and reported for each experiment.

The full replication package includes:
\begin{itemize}
\item Python source code for the model (\texttt{projects/did\_transformer/model/}), training (\texttt{train/}), synthetic validation (\texttt{synthetic/}), and analysis (\texttt{analysis/}).
\item R scripts for traditional DiD (\texttt{output/apep\_0489/v1/code/}).
\item Pre-trained model checkpoint and vocabulary files.
\item All synthetic DGP specifications and benchmark results (JSON).
\end{itemize}


\end{document}
